\documentclass[11pt]{article}
\usepackage[a4paper,margin=26mm]{geometry}
\usepackage[T1]{fontenc}
\usepackage{amsmath,amssymb,amsthm,booktabs,hyperref}
\hypersetup{colorlinks=true,urlcolor=blue,linkcolor=blue}
\newtheorem{theorem}{Theorem}
\newtheorem{lemma}[theorem]{Lemma}
\newtheorem{proposition}[theorem]{Proposition}
\newtheorem{corollary}[theorem]{Corollary}
\newcommand{\TV}{\mathrm{TV}}
\newcommand{\Unif}{\operatorname{Unif}}
\newcommand{\Law}{\mathcal L}
\usepackage{lmodern}


\title{A general convergence proof for protected/MIN CYCLON}
\author{Standalone research proof for independent review}
\date{16 November 2026}
\begin{document}
\maketitle
\begin{abstract}
For the precisely specified fixed-membership, atomic, independently
scheduled algorithm below, we prove a common atom for every
$1\le\ell\le c\le n-1$, and consequently a unique stationary full-state
law, uniform geometric total-variation convergence, and conditional
uniform local sampling at deterministic times. The rate is implicit,
not a polynomial or practical mixing bound. This is a working mathematical
proof, not a formally verified or externally peer-reviewed result.
Only the algorithm, assumptions, theorem and supporting proof chain
are included here. The shuffle-length-one, two and three arguments are
necessary induction bases, not numerical experiments.
\end{abstract}
\tableofcontents


\section{Algorithm and assumptions}
\label{sec:target}

Throughout this document the transition kernel is $\widehat P_\ell$
(also denoted $P$ when parameters are fixed). All statements use the
following single protected/MIN mandatory-removal algorithm.

\subsection{System, initialization, and randomness}
\begin{enumerate}
\item Membership is a fixed set $V=\{1,\ldots,n\}$ with $n\ge2$.
All peers follow the algorithm. There are no joins, departures,
crashes, message losses, or external introductions of identities.
\item Every peer has the same cache capacity $c$ and fixed shuffle
length $\ell$, with $1\le\ell\le c\le n-1$. A cache $C_v$
contains at most $c$ descriptors $u^a$, where $u\ne v$ and
$a\in\mathbb Z_{\ge0}$. No cache contains two descriptors for
the same identity. Empty and underfull caches are allowed.
\item Initial ages are arbitrary finite nonnegative integers,
without any common a priori bound. The initial directed overlay
has an edge $v\to u$ when $u\in C_v$; its underlying undirected
graph is connected. Strong connectivity and initially full caches
are not assumed.
\item At each global event choose an initiator independently and
uniformly from $V$, including peers with empty caches. Equivalently,
independent rate-one Poisson clocks at the peers generate this
embedded event chain. Positive results for this scheduler do not
automatically cover fixed periodic phases.
\item An exchange is reliable and atomic. Both buffers are formed
before either cache processes incoming entries. No other exchange
intervenes. Oldest-age ties, extra subsets, reply subsets, and
eligible victim sets are chosen uniformly. Conditional on the
current state and earlier choices in that exchange, the random
choices are independent, including victim choices at its two
endpoints. Fresh choices are used at subsequent events.
\end{enumerate}

\subsection{One activation and exchange}
For a packet or cache $A$, write $\operatorname{ids}(A)$ for
its identity set. Removing an identity removes its descriptor.
An event with initiator $p$ is the following procedure.
\begin{enumerate}
\item If $C_p$ is empty, do nothing and end the event.
Otherwise increment the age of every descriptor in $C_p$ by one.
\item Choose a maximum-age identity $q\in C_p$, uniformly among
ties. Choose a uniform subset $S$ of
$\min(\ell-1,|C_p|-1)$ other descriptors from $C_p\setminus\{q\}$.
The outgoing request packet is
\[
 B_p=\{p^0\}\cup S.
\]
Thus the partner's descriptor is replaced in the packet by a
fresh self descriptor, while the extra descriptors keep their
incremented ages.
\item Peer $q$ forms its reply $B_q$ as a uniform subset of
$\min(\ell,|C_q|)$ descriptors from its pre-exchange cache.
It does not increment its ages when acting as responder.
\item Remove the contacted partner $q$ unconditionally and
apply the merge procedure below at both endpoints, using the
already formed packets:
\[
\begin{split}
 C_p'&=\operatorname{Merge}_p
       (C_p\setminus\{q\},B_q,\operatorname{ids}(S)),\\
 C_q'&=\operatorname{Merge}_q
       (C_q,B_p,\operatorname{ids}(B_q)).
\end{split}
\]
Here the right-hand $C_p$ has the incremented ages from step~1;
the displayed set difference performs the mandatory removal.
All other caches
remain unchanged.
\end{enumerate}

\paragraph{Merge procedure.}
$\operatorname{Merge}_v(A,R,E)$ receives a resident cache $A$,
an incoming packet $R$, and the identity set $E$ of sent entries
eligible for replacement. Compute
\[
 D=\operatorname{ids}(R)\setminus\{v\},\qquad
 N=D\setminus\operatorname{ids}(A),\qquad
 d=\bigl(|N|-(c-|A|)\bigr)_+ .
\]
Ignore the incoming self descriptor, if any. For every identity
in $D\cap\operatorname{ids}(A)$ retain one descriptor with
age equal to the minimum of its resident and incoming ages.
These received identities are protected from eviction.
Choose a uniform $d$-element subset of
$E\cap\operatorname{ids}(A)\setminus D$ as victims, delete
those residents, and insert all descriptors with identities
in $N$, carrying their received ages. Newly inserted entries
cannot be victims, and unsent residents cannot be victims.
Lemma~\ref{lem:receiveprotect} proves that enough eligible
victims always exist, including for underfull exchanges.

The two key rules are therefore explicit: a valid received
identity is present after merging, and a duplicate keeps the
younger age. Sending an extra descriptor alone does not delete
its resident copy; it is deleted only if selected as an eligible
victim. The contacted partner is removed unconditionally.




\subsection{Notation, observation rule and proof organization}
For an ordered population $V=\{v_1<\cdots<v_n\}$, the canonical state
$H(V,c)$ gives $v_i$ the $\min(c,n-i)$ greatest labels above it, all at
age zero. The greatest vertex (the root) is empty; every other vertex
caches it. For singleton components $H$ is the empty state. In a closed
subsystem with $m$ identities the effective parameters are
$\min(c,m-1)$ and $\min(\ell,m-1)$; singleton components need no events.
A \emph{universal zero-age anchor} $a$ is an identity stored as $a^0$
by every nonowner; its owner cache is unrestricted. A \emph{holder}
of $a$ is a nonowner caching $a$; a \emph{nonholder} always excludes
the owner. Write $N_a=|\{v\ne a:a\in C_v\}|$ and
$Z_a=|\{v\ne a:a^0\in C_v\}|$. Caches used in set expressions mean
identity sets unless ages are displayed.

All controlled paths below prescribe possible outcomes of the stated
random choices; they do not alter the algorithm or require peers to
know the global state. A finite legal path has positive probability.
Uniform bounds require additional argument, since ages are unbounded.
At a deterministic event time an observer $i$ requesting $1\le b\le c$
samples uniformly without replacement from its cache and orders the
sample uniformly, conditional on $|C_i|\ge b$. This does not assert
always-available samples, independence across peers or times, uniformity
of whole overlay graphs, or correctness for stopping-time queries,
churn, periodic schedules, or other replacement conventions.

The final general theorem is Theorem~\ref{thm:allparametersuniform}.
Sections below first establish general tools and the three induction
bases, then the lower-length structural interface, the independent
full-capacity theorem, and finally the arbitrary-capacity source proof.


\subsection{Main theorem}

\begin{theorem}[General fixed-length convergence for the adopted kernel]
\label{thm:allparametersuniform}
For every $n\ge2$ and $1\le L\le c\le n-1$, the adopted
protected/MIN mandatory-removal kernel on weakly connected states
with arbitrary finite nonnegative integer ages has a uniform
finite-horizon common atom at any prescribed labelled $H(V,c)$.
Consequently it has a unique stationary full-state law and converges
to it in total variation, uniformly at an implicit geometric rate.
Every allowed sample size has positive stationary availability, and
conditional ordered local samples are uniform in stationarity and
converge to that law at deterministic observation times.
No polynomial or practically useful mixing bound is asserted.
\end{theorem}

The proof appears in the final section, after its supporting lemmas.

\section{Connectivity and probabilistic tools}

\begin{lemma}[Protection is feasible and preserves connectivity]
\label{lem:receiveprotect}
Received-identity protection is feasible for every $n,c,\ell$
and every underfull exchange. Under this rule every weak
component is preserved.
\end{lemma}
\begin{proof}
At the initiator let $v$ be its vacancy count after mandatory
partner removal, $k$ the number of extra sent residents, and
$e$ the number of those residents protected by received
identities. If $i$ incoming entries are usable, then
$i+e\le\ell$ and $v+k\ge\ell$.
Thus $(i-v)_+\le k-e$, so enough unprotected eligible
victims remain. At the responder use its vacancy count $v$,
reply size $m$, protected eligible count $e$, and usable
incoming count $j$. Again $j+e\le\ell$ and $v+m\ge\ell$,
so $(j-v)_+\le m-e$. Newly inserted entries are not victims.

The responder retains or inserts the initiator's identity,
which is always in the outgoing packet; hence the removed
partner edge is preserved in reverse.
Every removed initiator extra was sent and therefore is
retained or inserted at the responder.
Every removed responder entry was returned and is retained
or inserted at the initiator. The latter cannot be the
initiator's self identity, because it is protected at the
responder. Each lost edge has a path through the two
exchanging peers. New references stay within their previous
weak component.
\end{proof}

\begin{proposition}[Lossless exchange of identity knowledge]
\label{prop:thesisunion}
For every shuffle length and capacity, each exchange $p\to q$
under $\widehat P_\ell$ preserves
\[
 K_p'\cup K_q'=K_p\cup K_q,\qquad K_v=C_v\cup\{v\},
\]
where caches in this formula denote identity sets.
In particular every weak component is preserved.
\end{proposition}
\begin{proof}
Every removed nonendpoint identity was sent to the opposite
endpoint and is retained there by received-identity protection.
The endpoint identities remain in their own knowledge sets.
Conversely, every new cached identity came from one of the
two old caches or is an endpoint's self identity.
This proves union equality. The edge-transfer proof of
Lemma~\ref{lem:receiveprotect} gives weak connectivity.
Minimum-age merging changes neither assertion about identities.
\end{proof}

\subsection{From a common atom to conditional uniformity}

\begin{lemma}[A common atom gives convergence, including ages]
\label{lem:atom}
Let $P$ be a Markov kernel on a countable state space. Suppose there
are $L\ge1$, $\delta>0$, and a state $z$ such that
\[
 P^L(x,\{z\})\ge\delta\qquad\text{for every state }x.
\]
Then $P$ has a unique stationary probability law $\pi$, and
\[
 \sup_x\|P^t(x,\cdot)-\pi\|_{\TV}
 \le(1-\delta)^{\lfloor t/L\rfloor}.
\]
Irreducibility of the whole state space is not required.
\end{lemma}
\begin{proof}
The case $\delta=1$ is immediate. Otherwise split
$P^L(x,\cdot)=\delta\delta_z+(1-\delta)R(x,\cdot)$, where $R$ is a
Markov kernel. Hence $\mu\mapsto\mu P^L$ is a contraction by
$1-\delta$ in total variation. Probability measures on a countable
space are complete in this metric, so the contraction has a unique
fixed point $\pi$. Since
$(\pi P)P^L=(\pi P^L)P=\pi P$, uniqueness gives $\pi P=\pi$.
Every $P$-stationary law is also $P^L$-stationary, proving uniqueness
for $P$. Iterating the contraction and applying the remaining
$t\bmod L$ transitions proves the bound.
\end{proof}

\begin{theorem}[Stationary symmetry and availability]
\label{thm:stationary}
Let $P$ be equivariant on a relabelling-invariant state space.
Suppose it has a unique stationary law $\pi$ and
$\eta_t(x)=\|P^t(x,\cdot)-\pi\|_{\TV}\to0$.
Fix $i$ and $1\le b\le c$, and suppose
$a=\pi\{|C_i|\ge b\}>0$. For uniform sampling without replacement
from $C_i$, followed by uniform ordering,
\[
 \Law_\pi(S_i^{(b)}\mid |C_i|\ge b)
 =\Unif((V\setminus\{i\})_b).
\]
The same conditional sample law converges to this uniform law from
every initial state $x$. If $\eta_t(x)\le a/2$, its distance is at most
$4\eta_t(x)/a$.
\end{theorem}
\begin{proof}
For every permutation $\sigma$ of $V$, equivariance makes
$\sigma_*\pi$ stationary. Uniqueness gives $\sigma_*\pi=\pi$.
The event $A=\{|C_i|\ge b\}$ is invariant under permutations fixing
$i$, and the output rule is equivariant under this group. Its action
on distinct ordered $b$-tuples of other identities is transitive.
An invariant probability law on that finite set is uniform.

For the quantitative claim include the extra randomness of sample
extraction in the transition from a state to $(A,S_i^{(b)})$.
Total variation cannot increase under this transition. Put
$a_t=\Pr_x(A\text{ at }t)$. Then $|a_t-a|\le\eta_t(x)$.
For any set $B$ of available tuples,
\[
 \left|\frac{\Pr_x(A,S_i^{(b)}\in B)}{a_t}
       -\frac{\Pr_\pi(A,S_i^{(b)}\in B)}a\right|
 \le\frac{2\eta_t(x)}{a_t}
 \le\frac{4\eta_t(x)}a .
\]
Taking the supremum over $B$ proves the result.
\end{proof}

\begin{lemma}[Closed classes reached from a finite reset family]
\label{cor:thesisgeneralregimes}
Suppose a countable-state chain has a finite nonempty family $F$ and
constants $h\ge1$, $\delta>0$ with $P^h(x,F)\ge\delta$ for every
state. Then every $f\in F$ has a finite positive-probability path to
a closed communicating class. There are finitely many such classes.
If every such class reaches the same state $z$ and $P(z,\{z\})>0$,
then there are $T<\infty$, $\gamma>0$ such that
$P^T(x,\{z\})\ge\gamma$ for every $x$.
\end{lemma}
\begin{proof}
On $F$ put an edge $f\leadsto g$ whenever the full chain admits a
finite positive path from $f$ to $g$. Every vertex reaches a terminal
strongly connected component of this finite graph. Choose $z_k$ in
one such component, and let $C_k$ be all full-chain states reachable
from $z_k$. Any $y\in C_k$ admits a reset path to some $g\in F$.
Terminality puts $g$ in the same component as $z_k$, so $y$ returns
to $z_k$. Thus $C_k$ is communicating and closed. Distinct terminal
components give disjoint classes; every closed class intersects $F$
and is one of them. This proves the first assertions without assuming
irreducibility. Under the last hypothesis choose a path from every
$f\in F$ to $z$, pad these finitely many paths at $z$ to their maximum
length $m$, and let $\varepsilon>0$ be their minimum probability.
The reset followed by those paths gives $T=h+m$ and
$\gamma=\delta\varepsilon$.
\end{proof}

\section{Singleton shuffles: the first induction base}

\begin{theorem}[Finite-set regeneration for $\ell=1$]
\label{thm:agereset}
Use the algorithm of Section~\ref{sec:target}
with arbitrary $n,c$ and $\ell=1$. Label the peers $1,\ldots,n$ and set
\[
 K=n+c-1,\qquad L=nK,\qquad \delta=(nc)^{-L}.
\]
Let $\mathcal F$ be the finite set of states with all ages zero and
$C_i\subseteq\{i+1,\ldots,n\}$ for every $i$. Then, from every state
with arbitrary finite ages,
\[
 P^L(x,\mathcal F)\ge\delta.
\]
If $\tau_{\mathcal F}$ is its first hitting time, then
\[
 \Pr_x\{\tau_{\mathcal F}>mL\}\le(1-\delta)^m,\qquad
 \mathbb E_x\tau_{\mathcal F}\le L/\delta .
\]
\end{theorem}
\begin{proof}
First activate one fixed peer $p$ for $K$ consecutive events.
Only its cache ages are incremented during this block. When its
chosen partner $q$ already knows $p$, prescribe that $q$ return the
descriptor for $p$. This occurs with probability at least $1/c$.
The reply is a self entry at $p$, so mandatory removal reduces
$|C_p|$ by one. The request contains only $p^0$, a duplicate at
$q$; hence $q$ retains its descriptor for $p$.

If $q$ does not know $p$, its receipt of $p^0$ makes it learn $p$.
It either uses a vacancy or replaces its single returned entry.
Subsequent exchanges in this block cannot make it forget $p$:
it never initiates, and every request from $p$ contains only the
already-known identity $p$. There are at most $n-1$ such first-learning
events. The size of $C_p$ never increases, because one descriptor is
removed and at most one is received. Every other nonempty activation
in the prescribed block shrinks it. At most $c$ shrinkings empty it.
Consequently $n-1+c=K$ activations suffice, independently of all ages.
Remaining activations of an empty $p$ do nothing.

Now execute this block for $p=1,2,\ldots,n$ in order. Immediately
after its block, $p$ is empty. During later blocks it never initiates,
and every incoming request contains only a fresh self descriptor
$j^0$ from a later peer $j>p$. Thus its final cache consists solely
of age-zero descriptors of later peers. The final state belongs
to $\mathcal F$.

Prescribing the $L$ initiators has probability $n^{-L}$. Each reply
that must be prescribed has conditional probability at least $1/c$;
there are at most $L$ of them. No particular oldest-partner or
victim choice was required. The resulting lower bound is
$(nc)^{-L}$. Applying it to successive $L$-event blocks, using the
Markov property, gives the tail bound, and summing that bound gives
the expectation.
\end{proof}

\begin{theorem}[Every network size, capacity one]
\label{thm:mandatoryone}
Let $n\ge2$ and $c=\ell=1$ in
Section~\ref{sec:target}. On the weakly connected state
space, the full cache-and-age chain has a unique stationary law
$\pi$: the uniform law on all labelled rooted spanning trees,
with all descriptor ages zero and all arcs directed toward the
empty root. Set $D=2n(n-1)$ and $\beta=n^{-D}$. Then
\[
 \sup_x\|P^t(x,\cdot)-\pi\|_{\TV}
 \le(1-\beta)^{\lfloor t/D\rfloor}.
\]
At stationarity, every observer $i$ has
\[
 \pi(C_i=\varnothing)=\frac1n,\qquad
 \pi(C_i=\{j^0\})=\frac1n\quad(j\ne i).
\]
Thus the sample converges to uniform on the other peers conditional
on a nonempty cache, with limiting availability $1-1/n$.
\end{theorem}
\begin{proof}
Represent a singleton cache by its identity $f(p)$ and an empty
cache by $f(p)=\bot$. If $p$ is nonempty, put $q=f(p)$ and
$r=f(q)$. The identity update is
\[
 f'(q)=p,\qquad
 f'(p)=
 \begin{cases}
  \bot,& r\in\{p,\bot\},\\
  r,& r\notin\{p,\bot\}.
 \end{cases}
\]
Other pointers do not change. A path on three distinct vertices
$p\to q\to r$ becomes $q\to p\to r$. Contact with an empty
responder reverses the edge. A reciprocal exchange removes only
one of its two opposite arcs. Weak connectivity is therefore
preserved.

A weakly connected digraph of outdegree at most one consists
either of a directed cycle with trees feeding into it, or of a
rooted spanning tree with one empty root. Indeed, following outgoing
pointers reaches a cycle or an empty vertex, and different terminal
cycles or empty vertices cannot have intersecting weak basins.
On a directed cycle of length $k\ge3$, activating a cycle vertex
shortens that cycle by one. On a two-cycle, activating either
endpoint produces a rooted tree. At most $n-1$ prescribed
activations thus reach a rooted tree, which remains a rooted tree.

In a rooted tree the initiator's descriptor is deleted, the
responder's descriptor transfers to the initiator if present, and
the responder receives a fresh age-zero descriptor. There is no
copying, and the only incremented descriptor is deleted.
Consequently the number of positive-age descriptors decreases by
one whenever a positive-age carrier initiates, and never increases.
At most $n-1$ further prescribed activations clear every positive
age. The zero-age rooted-tree family is closed.

We next prove that its identity transition graph is connected.
Fix an anchor $a$ and follow its directed path to the root.
While an off-path vertex remains, choose one, $p$, whose parent
$q$ is on the path. Activating $p$ inserts $p$ immediately after
$q$ on the path; if $q$ is the root, $p$ becomes the new root.
Such a vertex exists by connectivity of the tree.
At most $n-1$ moves give a Hamiltonian path beginning at $a$.

On a Hamiltonian path with consecutive vertices $x,u,v$, write
$y=f(v)$, allowing $y=\bot$ when $v$ is the root.
The word $x,x,u,v$ swaps $u,v$ and changes no other outgoing edge:
\[
\begin{array}{c|ccc}
 &f(x)&f(u)&f(v)\\ \hline
 \text{initial}&u&v&y\\
 x&v&x&y\\
 x,x&y&x&x\\
 x,x,u&u&y&x\\
 x,x,u,v&v&y&u
\end{array}
\]
All four moves are valid, including when $y=\bot$.
Adjacent swaps sort the $n-1$ vertices following $a$, using at most
$(n-1)(n-2)/2$ swaps.
Every nontrivial tree move initiated by $p$ is inverted by
initiating its old parent $q$ in the resulting tree, including
when that parent was the empty root. Therefore we may reverse
the insertion path of a target tree. The tree transition graph
has diameter at most
\[
 2(n-1)+4\,\frac{(n-1)(n-2)}2=2(n-1)^2.
\]

Choose one fixed zero-age rooted star $z$.
The cycle reduction, age clearing, and tree diameter bound give
a path to $z$ of length at most
$2(n-1)+2(n-1)^2=D$ from every initial aged state.
Pad to exactly $D$ events using the empty root's no-op.
Every prescribed initiator has probability $1/n$; there are no
nontrivial buffer or tie choices at capacity one. Hence
$P^D(x,\{z\})\ge n^{-D}$.
Lemma~\ref{lem:atom} gives uniqueness and the displayed bound.

Pairing each marked tree transition $(f,p,g)$ with its inverse
$(g,q,f)$ proves symmetry of the identity kernel on the finite
tree family. Its uniform law, with all ages zero, is therefore
stationary and must equal $\pi$. Relabelling makes the root
uniform on $V$. Permutations fixing $i$ make the remaining
$n-1$ possible cached identities equiprobable, each with mass
$(1-1/n)/(n-1)=1/n$. Conditional convergence follows from
Theorem~\ref{thm:stationary}.
\end{proof}

\begin{theorem}[Arbitrary population and cache capacity, $\ell=1$]
\label{thm:generalone}
Let $n\ge2$ and $2\le c\le n-1$ in
Section~\ref{sec:target}, with $\ell=1$.
Define, using this fixed capacity $c$,
\[
 B_1=0,\qquad B_m=2m(m+c-1)+2B_{m-1}\quad(2\le m\le n),
\]
\[
 T=n(n+c-1)+B_n,\qquad \gamma=(nc^2)^{-T}.
\]
From every weakly connected state with arbitrary finite integer ages,
\[
 P^T(x,\{H(V,c)\})\ge\gamma .
\]
The full cache-and-age chain therefore has a unique stationary
law $\pi$ and
\[
 \sup_x\|P^t(x,\cdot)-\pi\|_{\TV}
 \le(1-\gamma)^{\lfloor t/T\rfloor}.
\]
For every observer and every $1\le b\le c$, availability has a
positive stationary probability, and the local ordered $b$-sample
converges to uniform conditional on availability.
Together with Theorem~\ref{thm:mandatoryone}, this covers every
allowed capacity for shuffle length one.
\end{theorem}

\subsection{Clearing blocks and separating one source}

Write $\mathcal B(p)$ for the clearing block in
Theorem~\ref{thm:agereset}: activate only $p$; whenever its partner
already knows $p$, prescribe reply $p$, and otherwise choose any
reply. With $m$ active vertices it empties $p$ within $m+c-1$
events. No unselected resident of the initiator can be replaced
when $\ell=1$, since the sole usable reply fits the vacancy made
by partner removal. If all ages are zero at the block's start,
they are zero at its end: only $p$ increments ages, and its cache
is empty at the end; passive peers receive only fresh $p^0$.

\begin{lemma}[An unknown source with a connected remainder]
\label{lem:source}
From any zero-age weakly connected state on $V$, clear its
smallest vertex $a$, then clear each vertex of
$W=V\setminus\{a\}$, keeping $a$ passive. At the end all ages
are zero, no cache in $W$ contains $a$, $C_a$ is nonempty, and
the graph induced by $W$ is weakly connected.
\end{lemma}
\begin{proof}
Immediately after clearing $a$, every component of the induced
graph on $W$ has a descriptor for $a$: $C_a$ is empty while the
whole graph is connected.
During the remaining blocks the component partition on $W$
can only coarsen. A $W$-internal exchange preserves its internal
connectivity by edge reversal and transfer; if its reply is $a$,
it deletes no $W$ edge. An exchange with passive $a$ likewise
deletes no $W$ edge and may add one through the reply.

The first contact with $a$ inserts its first cached identity and
defines an anchor component in $W$. Thereafter every cached
identity at $a$ belongs to that anchor, which may grow.
Indeed a new initiator $p$ contacting $a$ receives an anchor
identity, joining the anchor, or receives itself and was already
in it. A nonempty cache at passive $a$ cannot become empty:
inserting a new initiator either fills a vacancy or replaces one
entry, while a duplicate changes nothing.

Every initial component joins the anchor. A descriptor for $a$
can disappear either when its holder contacts $a$, joining the
anchor directly, or when a full responder returns it and replaces
it by the active peer's fresh identity. In the latter case it
transfers to, or is already present at, that active peer. The
reversed edge links the two within $W$, and its clearing block
must eventually contact $a$. Thus eliminating a component's
last such descriptor forces that component into the anchor.

At the end no $W$ cache contains $a$: each vertex was emptied
and later passive requests contain only other initiators'
identities. Consequently all components have joined, and at
least one contact has made $C_a$ nonempty. Since $a$ is now
globally unknown, it cannot reappear while $a$ stays inactive.
All block boundaries have zero ages.
\end{proof}

\subsection{Extending the canonical remainder}

\begin{lemma}[Setting the cache of an unknown source]
\label{lem:extend}
Let $c\ge2$. Suppose all ages are zero, the smallest vertex $a$
is unknown in $W=V\setminus\{a\}$, the induced state on $W$
is $H(W,c)$, and $C_a$ is any nonempty cache in $W$.
There is a sequence of $|V|$ clearing blocks ending with all
ages zero, $a$ still unknown in $W$, $W$ weakly connected,
and $C_a$ equal to the $\min(c,|W|)$ greatest members of $W$.
\end{lemma}
\begin{proof}
Let $m=\max W$, $d=\min(c,|W|)$, let $J$ be the $d$
greatest members of $W$, and put $S=J\setminus\{m\}$.

First clear $a$, using reply $a$ from a responder that knows
$a$, and otherwise its smallest cached identity.
Every $W$ vertex other than $m$ retains its edge to $m$.
A full responder has at least two entries, so its smallest
entry is not $m$; an underfull responder evicts nothing.
The block must visit $m$. Otherwise let $q$ be the greatest
visited vertex. At its first visit its untouched canonical cache
returns a label greater than $q$. That label either is already
pending at $a$ or is inserted, and cannot disappear without being
visited, contradicting maximality of $q$.
Since $m$ initially is empty and remains passive, this phase
ends with $C_m=\{a^0\}$ and $C_a=\varnothing$.
The remainder stays connected through its edges to $m$.

Next, for each $j\in S$, start its clearing block with $j\to m$
and prescribe reply $a$. Preserve the descriptor for $m$ at every
unprocessed vertex: at a full responder not knowing $j$, choose
any reply other than $m$; at one knowing $j$, reply $j$.
Underfull responders need not evict anything.
Thus each unprocessed $j$ still knows $m$ when its block begins.
All ages are zero, so its first choice of $m$ is a legal tie.
Before that contact $C_m$ consists of $a$ and the previously
processed members of $S$, with size at most $d-1<c$.
Fresh $j$ therefore fills a vacancy and preserves $a$ there.
On later visits, $m$ already knows $j$ and replies $j$.
The active $j$ obtains or retains a descriptor for $a$, so must
contact $a$ before its block ends.
Passive $a$ has fewer than $c$ entries during this phase.
After these blocks,
\[
 C_a=S,\qquad C_m=\{a\}\cup S,
\]
all at age zero, and every processed $j$ is free of $a$.

Now clear all $h\in W\setminus J$.
Keep $a$ at passive $m$: if $m$ is full and does not know $h$,
choose a reply other than $a$, possible because $c\ge2$;
otherwise use the ordinary clearing policy.
At passive $a$, protect every member of $S$.
If this phase is nonempty, $|S|=c-1$, so there is at most one
unprotected slot. Fill a vacancy for a new $h$, or return and
replace the unique unprotected entry. For a duplicate $h$,
reply $h$. These choices respect the clearing-block rule.
At the end the only remaining descriptor for $a$ is at $m$:
every other vertex was cleared after its last possible
acquisition of $a$, and later fresh requests cannot reintroduce
that identity. Cache $C_a$ consists of $S$ and at most one
unprotected entry.

Finally start $m$'s clearing block with $m\to a$, a zero-age tie.
At $a$, reply with and replace its unprotected entry, if present;
otherwise fill its vacancy. Fresh $m$ is new there, since $m$
has never initiated in these phases. Hence $C_a$ becomes
exactly $J$. This removes the globally last descriptor for $a$.
Finish clearing $m$; it cannot contact $a$ again, so $C_a$ stays
fixed. All ages are zero at the end.

After the first phase $W$ was connected. Thereafter $a$ is
passive; exchanges with it only add $W$ edges, while internal
exchanges preserve $W$ connectivity. Thus $W$ remains connected.
The number of clearing blocks is
$1+|S|+|W\setminus J|+1=|V|$.
\end{proof}

\subsection{Induction and the convergence bound}

\begin{proof}[Proof of Theorem~\ref{thm:generalone}]
For the induction allow a fixed configured capacity $c\ge2$
even on a smaller active population. Excess cache slots cause
no difficulty; the cache is still a set of nonself identities
from that population. The one-vertex case is empty.

On a zero-age connected $m$-vertex state, apply
Lemma~\ref{lem:source}, then freeze its unknown smallest vertex
$a$. Use induction to make $W$ canonical. Nothing can contact
$a$, so its cache stays fixed. Apply Lemma~\ref{lem:extend} to
set $C_a$ to the greatest $\min(c,|W|)$ labels, leaving $a$
unknown, $W$ connected, and all ages zero.
Freeze $a$ again and make $W$ canonical a second time.
This reaches exactly $H(V,c)$.
The two lemmas use $m$ clearing blocks each, of at most
$m+c-1$ events, and the two recursive calls cost at most
$2B_{m-1}$. Thus the stated recurrence bounds the entire path.

From arbitrary ages first apply the clearing sweep of
Theorem~\ref{thm:agereset}, costing $n(n+c-1)$ events and
leaving a connected zero-age state. The total path length is
at most $T$. Pad shorter paths at the empty largest vertex of
$H(V,c)$. At each prescribed event the initiator has probability
$1/n$, the chosen oldest tie has probability at least $1/c$,
and the prescribed nonempty reply has probability at least $1/c$.
Empty replies and no-ops require fewer choices.
Therefore the path probability is at least $(nc^2)^{-T}$.
Lemma~\ref{lem:atom} proves the full-state conclusion.

For any observer, relabel the same construction so that it is
the smallest vertex of the canonical state and therefore has
$c$ entries. Stationarity and the same minorization imply that
its full-cache probability is at least $\gamma$.
The transition rule and connected state space are equivariant,
so Theorem~\ref{thm:stationary} proves conditional uniformity
for every $1\le b\le c$.
\end{proof}



\begin{corollary}[Singleton common atoms for the adopted kernel]
\label{thm:thesisone}
At length one, every allowed population and capacity has a uniform
finite-horizon common atom at any labelled $H(V,c)$.
\end{corollary}
\begin{proof}
The preceding proofs were stated directly for the protected/MIN
kernel. To make explicit where age merging matters, the clearing
block is partner-agnostic: repeat contacts return the initiator's
identity and shrink its row even if the passive copy has been
refreshed. Each clearing block starting at zero ends at zero.
Source separation and extension use only identities, cardinalities,
and zero-age ties at block boundaries. At capacity one, the
zero-age rooted-tree kernel is unchanged by MIN. Its connected tree
transition graph reaches every zero-age rooted star, which is
$H(V,1)$, and the empty root permits uniform padding.
For $c\ge2$ apply Theorem~\ref{thm:generalone}.
\end{proof}

\section{A finite-age reset at every shuffle length}

\begin{lemma}[A decreasing color potential for every $\ell\ge2$]
\label{lem:generalcolor}
Color descriptors young or old, color fresh self descriptors
young, and let mixed duplicate merges keep young. Assume every
peer holding an old descriptor has an old maximum-age descriptor,
so an old oldest partner can be chosen. Then, whenever any old
descriptor remains, legal controlled choices strictly decrease
lexicographically the nonnegative integer triple
\[
 \Phi=\left(nc-F,\ V,\ W\right),\qquad
 F=\sum_i f_i,\quad V=\sum_i f_i^2,\quad
 W=\sum_i o_i\ell^{c-f_i},
\]
where $f_i$ and $o_i$ are the young and old cache counts at
peer $i$. The coordinates lie in the ranges
\[
 0\le nc-F\le nc,\quad 0\le V\le nc^2,\quad
 0\le W\le nc\ell^c.
\]
\end{lemma}
\begin{proof}
Choose an initiator $p$ maximizing $f_p$ among peers with
$o_p>0$, and choose a legal old partner $q$.
Write $P,Q$ for their pre-update identity sets, $Y_p,Y_q$
for their young identity sets, $a=|P|$, $b=|Q|$, and
\[
 k=\min(\ell-1,a-1),\qquad m=\min(\ell,b).
\]
Choose the extra set $S\subseteq P\setminus\{q\}$ of size
$k$, including one member of $Y_p\setminus Y_q$ if that set
is nonempty. Such an extra exists when needed, since $q$
is old and $\ell\ge2$. Put $S_Y=S\cap Y_p$,
$S_O=S\setminus S_Y$, and $U=\{p\}\cup S$.
The fresh self identity $p$ is young in the outgoing packet.

The responder needs
\[
 r_q=\bigl(b+|U\setminus Q|-c\bigr)_+
\]
victims. Forbid victims in
\[
 E=(Q\cap U)\cup\bigl(Y_q\cap(P\cup\{p\})\bigr).
\]
The first part is required by received-identity protection;
the second prevents loss of a young copy already young at $p$.
If $d_0=|Q\cap U|$, then
\[
 |E|\le d_0+(a-1-k),\qquad
 b-|E|\ge b-d_0-a+1+k.
\]
The extra term counts only unsent identities of $P$, since
shared sent identities are already in $Q\cap U$, and
$q\notin Q$. When $r_q>0$,
$r_q=b+k+1-d_0-c$, so $b-|E|\ge r_q$ by $a\le c$.
Choose $r_q$ victims $D_q\subseteq Q\setminus E$.

Require the reply $T\subseteq Q$ to contain
\[
 D_q\cup A,\qquad
 A=(S_Y\cap Q)\cup(\{p\}\cap Q).
\]
These two sets are disjoint and $A\subseteq Q\cap U$.
If $r_q>0$, their total size is at most
$r_q+d_0=b+k+1-c\le\min(\ell,b)=m$.
If $r_q=0$, $|A|\le d_0\le m$ instead.
Fill the reply arbitrarily to size $m$.
Every chosen responder victim was sent and is unprotected,
so this is a legal joint outcome of reply and victim choices.
At the initiator, choose old eligible victims before young
ones, using only the already permitted victim sets.

We first prove that $F$ cannot decrease, with strict increase
if an outgoing young identity was not already young at $q$.
Let
\[
 d_Y=|S_Y\cap Q|,\quad \sigma=\mathbf1_{\{p\in Q\}},\quad
 e_O=|S_O\cap T|.
\]
At most $m-d_Y-\sigma-e_O$ reply identities are usable
new entries at $p$: the counted entries are duplicates or
the self identity. There are $c-a+1$ mandatory vacancies
and $k-|S_Y|-e_O$ eligible old extras. Old-first victim
selection therefore loses at most
\[
 \bigl(|S_Y|-d_Y-\sigma\bigr)_+=(j_Y-1)_+,
 \qquad
 j_Y=\mathbf1_{\{p\notin Q\}}+|S_Y\setminus Q|,
\]
young copies at $p$. Here we used
$c-a+1+k\ge\ell\ge m$.
Every young victim at $q$ is absent from $P\cup\{p\}$,
so its reply inserts a new protected young copy at $p$.
These young losses are compensated one for one.
The outgoing packet additionally creates $j_Y$ young copies
at previously absent identities, and converts $g$ old copies
at $q$ to young, where $g$ counts its outgoing young
identities previously old there. Consequently
\[
 F'-F\ge j_Y+g-(j_Y-1)_+\ge0,
\]
with strict inequality if $j_Y+g>0$.

Suppose now that $F'=F$. Then $p\in Y_q$ and every
young member of $S$ belongs to $Y_q$. By the prescribed
choice of $S$, this implies $Y_p\subseteq Y_q$, hence
\[
 f_q\ge f_p+1.
\]
In particular $q$ has no old descriptor: otherwise its
young count would contradict the maximal choice of $p$.
No young copy is lost at $p$. All young changes consist
of $t$ responder victims transferring to new young locations
at $p$, so
\[
 f_p'=f_p+t,\quad f_q'=f_q-t,\quad
 0\le t\le f_q-f_p-1.
\]
For $t>0$ this gives
\[
 V'-V=2t\bigl(t-(f_q-f_p)\bigr)<0.
\]

If also $V'=V$, then $t=0$. The responder was entirely
young and no new young copy is created at $p$, so $p$
receives no new descriptor at all. It loses exactly its
selected old partner. The responder can acquire at most
$\ell-1$ old extras, while the young counts stay unchanged.
Therefore
\[
 W'-W\le-\ell^{c-f_p}
       +(\ell-1)\ell^{c-f_q}<0,
\]
using $f_q\ge f_p+1$.
This proves strict lexicographic decrease in every case.
The coordinate bounds follow directly from capacity $c$.
\end{proof}

\begin{theorem}[Uniform finite-age reset for every fixed shuffle length]
\label{thm:thesisgeneralreset}
Use the fixed adopted algorithm in Section~\ref{sec:target},
with $2\le\ell\le c\le n-1$. Set
\[
 H=(nc+1)(nc^2+1)(nc\ell^c+1),\qquad
 B=(nc+1)H,\qquad
 \delta=(nc\,16^c)^{-H}.
\]
For the finite family $F_B$ of weakly connected aged states
whose maximum age is at most $B$,
\[
 \widehat P_\ell^H(x,F_B)\ge\delta
\]
uniformly over all weakly connected initial states with
arbitrary finite integer ages. In particular,
\[
 \Pr_x\{\tau_{F_B}>kH\}\le(1-\delta)^k,\qquad
 \mathbb E_x\tau_{F_B}\le H/\delta.
\]
The construction does not require occasional shorter shuffles.
\end{theorem}
\begin{proof}
Among the $nc+1$ intervals $(jH,(j+1)H]$, $j=0,\ldots,nc$,
one contains no initial age, since there are at most $nc$
initial descriptors. Put $a=jH$ for that interval.
Color initial ages at most $a$ young and those greater
than $a+H$ old. Fresh self descriptors are young.
For the first $H$ events, young ages are at most $a+t\le a+H$
and old ages remain strictly greater than $a+H$, even after
old--old minimum merging. A mixed merge keeps the young copy.
Thus every oldest partner of a holder of an old descriptor
is old during this horizon.

Apply Lemma~\ref{lem:generalcolor} at every event until
old descriptors are gone. Its integer triple takes at most
$(nc+1)(nc^2+1)(nc\ell^c+1)=H$ values, and strictly decreases
while any old remains. Hence old-color extinction occurs
within $H$ events. Padding a shorter path by arbitrary
valid events cannot recreate old colors. The maximum age
at event $H$ is therefore at most $a+H\le B$.
Connectivity holds by Proposition~\ref{prop:thesisunion}.

Each prescribed event has probability at least
$(nc\,16^c)^{-1}$. A conservative product counts initiator
probability $1/n$, oldest-tie probability at least $1/c$,
the extra and reply subsets, and both eligible victim subsets.
The latter four subset denominators have product at most
$2^{c-1}2^c2^{\ell-1}2^\ell\le16^c$.
Prescribing a compatible reply/victim pair in the proof
does not change their actual order of sampling.
Multiplication along the controlled path gives $\delta$;
padding after extinction is unrestricted.
The Markov property on successive $H$-event blocks proves
the hitting bounds.
\end{proof}

\section{Two-entry shuffles: the second induction base}

\begin{lemma}[Removing a universal anchor from controlled paths]
\label{lem:anchorreduction}
Let $2\le\ell\le c$ and suppose $a$ is a universal
zero-age anchor. Freeze $a$ and remove its descriptor from
every other cache. Every finite path of the adopted algorithm
on $W=V\setminus\{a\}$ with capacity $c-1$ and shuffle
length $\ell-1$ can be simulated by a positive-probability
path of the original algorithm, preserving every anchor
copy at age zero and leaving $C_a$ unchanged. Reduced
activations of empty caches are omitted.
\end{lemma}
\begin{proof}
For a nonempty reduced cache at $p$, choose its prescribed
oldest residual partner $q$. This is a legal actual oldest
choice: after the increment every residual age is at least
one, while the anchor's age is exactly one.
Include $a$ among the outgoing extras and require $a^0$
in the reply. The remaining buffer sizes are exactly
\[
 \min(\ell-2,|C_p\setminus\{a\}|-1),\qquad
 \min(\ell-1,|C_q\setminus\{a\}|),
\]
as required for the reduced exchange. At $p$ the returned
$a^0$ refreshes its incremented anchor; at $q$, its resident
$a^0$ survives the incoming $a^1$ by minimum merging.
The anchor is protected at both endpoints. The other
duplicate merges, vacancy counts, and eligible victim sets
are exactly those of capacity $c-1$. All prescribed subset
and victim choices have positive probability.
If a reduced cache is empty, its reduced activation is a
no-op and is skipped, rather than activating its actual
singleton cache toward $a$. Thus this is a path simulation,
not an equality of the two randomly scheduled kernels.
\end{proof}

\begin{proposition}[A universal anchor suffices when $\ell=2$]
\label{prop:anchorcanonical}
For the adopted algorithm with $\ell=2$ and any $c\ge2$,
every state with a universal zero-age anchor reaches a
zero-age canonical state $H(V,c)$ by a finite
positive-probability path. The order defining $H$ may be
chosen arbitrarily. In particular all such labelled
canonical states communicate.
\end{proposition}
\begin{proof}
First keep the anchor $a$ fixed. By
Lemma~\ref{lem:anchorreduction}, the proved singleton
constructions can canonicalize each residual weak component
on $W$, with capacity $c-1$ and any chosen component root.
For capacity one use Theorem~\ref{thm:mandatoryone}; for
higher capacities use Theorem~\ref{thm:thesisone} and its
constructive proof, which permits excess slots on smaller
active populations. Singleton components need no exchanges.

We first empty $C_a$, without any assumption on its ages.
If it is nonempty, choose an oldest $q\in C_a$ and
canonicalize $q$'s residual component with $q$ as root,
while freezing $a$. Then $C_q=\{a^0\}$ and $C_a$ is
unchanged. Activate $a$ toward $q$. Its only reply is the
self identity $a$, so $C_a$ loses $q$ and receives nothing.
At $q$, fresh $a^0$ is protected and at most one extra
descriptor is inserted, which fits since $c\ge2$.
Every other anchor stays at age zero. Repeat at most $c$
times to empty $a$.

Now canonicalize each residual component at age zero.
Roots $u,v$ of two different components have
$C_u=C_v=\{a^0\}$. The following four exchanges merge
their residual components and restore the empty owner and
all anchor copies:
\[
\begin{array}{c|ccc}
 &C_a&C_u&C_v\\\hline
 \text{initial}&\varnothing&\{a^0\}&\{a^0\}\\
 u\to a&\{u^0\}&\varnothing&\{a^0\}\\
 v\to a&\{u^0,v^0\}&\varnothing&\{u^0\}\\
 a\to v&\{u^0\}&\varnothing&\{u^0,a^0\}\\
 a\to u&\varnothing&\{a^0\}&\{u^0,a^0\}
\end{array}
\]
At the third exchange choose $v$ in the oldest tie.
Its reply $u^0$ refreshes the initiator's $u^1$, and
the responder keeps $u^0$ when receiving $u^1$.
All displayed caches have at most two entries, so the
table is legal for every $c\ge2$; other peers are untouched.
The residual edge $v\to u$ merges the components.
Repeat with the same root $u$ until $W$ is connected,
then use the reduced singleton construction to obtain
$H(W,c-1)$. Together with the empty anchor $a$, this is
exactly $H(V,c)$ with $a$ ordered last.

To change the anchor to any desired $b\ne a$, first
canonicalize $W$ with $b$ last. Then $C_b=\{a^0\}$,
and every other peer in $W$ contains $b^0$.
One exchange $b\to a$, with empty reply, gives
$C_b=\varnothing$ and $C_a=\{b^0\}$ and leaves all
other $b^0$ copies unchanged. The result is the canonical
state with the final two labels $b,a$ interchanged.
With $b$ fixed as anchor, the reduced singleton
construction sets all remaining labels in any desired
order. This proves the assertion for every order.
\end{proof}

\begin{theorem}[A uniform first-contact bound for $\ell=2$]
\label{thm:firstcontact}
Use the adopted algorithm with $\ell=2$.
Fix a nonempty proper set $G\subset V$, put $S=V\setminus G$
and $s=|S|$, and suppose some cache in $S$ contains an
identity $r\in G$. Define
\[
 \begin{split}
 D_s&=s(s+c),\qquad J_s=(sc+1)(2s+c),\\
 T_s&=D_s(J_s+1)+J_s,\qquad
 \eta_s=(nc^5)^{-T_s}.
 \end{split}
\]
There is a positive-probability path, using only initiators
in $S$, that contacts $G$ within $T_s$ events.
Under the original uniform scheduler, if $\tau$ is the
first exchange between $S$ and $G$ in either direction,
\[
 \Pr_x\{\tau\le T_s\}\ge\eta_s,\qquad
 \Pr_x\{\tau>kT_s\}\le(1-\eta_s)^k.
\]
No bound on initial ages or on the frozen caches is needed.
The constructed path always activates an $r$ holder,
sends $r$ as extra when its partner is not $r$, and
requires $r$ in the reply whenever the responder has it.
The total number of $r$ holders cannot decrease along
this path except at its final contact if that partner
is owner $r$.
\end{theorem}
\begin{proof}
We construct a path while keeping $G$ inactive; any
contact with $G$ terminates the construction successfully.
Before such contact, identity $r$ remains in at least one
$S$ cache by received-identity protection and the pair-union
identity in Proposition~\ref{prop:thesisunion}.
No fresh $r$ can be generated inside $S$.
Its minimum outside age $m$ is therefore nondecreasing.

First we show that $m$ can be raised by at least one,
or $G$ contacted, within $D_s$ events. Pick a holder $p$
of $r^m$ and activate only $p$. Until success or loss of
this minimum copy, send $r$ as extra. If its actual oldest
partner is in $G$, we have succeeded. If a responder $q$
does not store $r^m$, require $r$ in the reply if it is
known there. The increment at $p$ destroys that
minimum copy and none is created: the outgoing copy has
age $m+1$, and the reply has no $r^m$.
If $q$ does store $r^m$, require it to return this
descriptor, plus $p$ when known or one other residual
entry otherwise. The initiator's copy is refreshed to $m$;
the responder's is protected and stays at $m$.
There is at most one usable new reply, so $p$ has no
extra eviction. Every such responder learns and thereafter
retains $p$. A repeat returns $r^m,p$, shrinking $C_p$.
Thus within $s+c$ events either a minimum copy is lost,
or $p$ is left with the singleton $r^m$ and contacts $G$.
There are at most $s$ minimum copies. Repeating this
operation proves the bound $D_s=s(s+c)$.

Repeat the minimum-raising construction $J_s+1$ times.
Unless $G$ has already been contacted, every outside
$r$ copy now has age greater than $J_s$.
Mark all descriptors currently present auxiliary old,
and all subsequent fresh self descriptors auxiliary young.
For this proof-only marking, a duplicate is young whenever
\emph{either} input is young, regardless of which supplies
the minimum numerical age. This marking is sound: the
minimum age is at most the young input's age, so after
$t$ further events every young descriptor has age at most
$t$. No $r$ descriptor becomes young, since its owner is
frozen. For the next $J_s$ events every selected oldest
partner of an $r$ holder is therefore auxiliary old.

The following pair-retaining routing policy now applies until
contact with $G$. Follow an $r$ holder $p$ and send $r$ as its extra
when its partner is $q\ne r$. Write $P=C_p$, $Q=C_q$. Both $r$ and
the selected partner are old. A responder entry is eligible for the
following victim choice if it is old, or young with identity outside
$P\cup\{p\}$. Exclude received identities $p,r$ from victims.
If $Q$ knows both $p,r$, reply with them. If it knows exactly one,
return that one and, if needed, one eligible victim different from
it. At a full responder there are $c-1$ other entries and at most
$|P|-2\le c-2$ forbidden young entries, so a victim exists.
If $Q$ knows neither, insert both, choosing
$(2-(c-|Q|))_+$ eligible victims and including them in the reply.
At most $|P|-2$ entries are forbidden young ones, leaving at least
$|Q|-c+2$ candidates. Complete replies to their prescribed size
arbitrarily. All these victims are sent and unprotected.

The responder ends knowing $p,r$. Each young victim moves to a
new young location at $p$, compensating its loss. At $p$ only the
old extra $r$ can be evicted, and the mandatory partner is old.
Young-OR MIN merging cannot lose a young marking. Hence the number
of young locations in $S$ never decreases. If $p$ loses $r$, both
reply entries must be new and the responder knew neither $p$ nor
$r$. Its new fresh $p$ gives a strict young-count increase; follow
$r$ at that responder. There are at most $sc$ such migrations.

Between migrations every visited responder learns and retains the
pair $p,r$. Repeat visits return that pair, both self or duplicate
at $p$, shrinking its cache. First visits can grow that cache by at
most one; there are at most $s-1$ nonterminal partners. Thus each
fixed-holder segment, including its terminal contact or migration,
has at most $2s+c$ events. The $(sc+1)$ segments contact $G$ within
$J_s=(sc+1)(2s+c)$ events. In particular the bounded-age young
marking remains valid for the entire routing stage.
Together with minimum raising this proves the bound $T_s$.

Each prescribed event has probability at least
$(nc^5)^{-1}$: the initiator contributes $1/n$, the oldest tie
and extra choice each at least $1/c$, the reply subset at least
$1/c^2$, and the only possibly nontrivial responder victim choice
at least $1/c$ (the active extra is unique). The denominator
is $n$, even though the witness path uses only $S$.
The holder-count claim follows directly from the packets:
an existing responder copy is protected, and its reply
protects the initiator's copy; a new responder copy
compensates any initiator loss. Only owner $r$ cannot
retain or return its own identity at the terminal contact.
This proves the one-block bound. Before any cross exchange,
$G$-internal events cannot change the outside caches, and
$S$-internal events preserve some outside $r$ copy.
The same hypothesis therefore holds at every block start
on $\{\tau>kT_s\}$. The Markov property gives the tail bound.
\end{proof}

\begin{lemma}[Source-only gain of a fresh copy]
\label{lem:c2sourcegain}
Fix $a$, and let $I=\{v\ne a:a^0\in C_v\}$.
If a vertex outside $I\cup\{a\}$ is reachable from $a$
by a directed cache path, a positive-probability path of
at most $(n-1)(n+1)$ activations of $a$ increases $|I|$.
No previously existing $a^0$ copy is lost.
\end{lemma}
\begin{proof}
Every source-only exchange delivers protected fresh $a^0$
to the responder, so $I$ cannot lose a member. Stop if a
responder is outside $I$, since that immediately gives a gain.
Until then every responder stores $a^0$ and returns it;
at most one other reply identity is usable at $a$.
It fits the mandatory vacancy, so an extra retained at
the initiator cannot be evicted.

Choose a shortest directed path from $a$ to outside
$I\cup\{a\}$, with first head $z$. Its internal vertices
belong to $I$. While contacting any other partner, send
$z$ as extra. An off-path responder learns and retains
$z$; if revisited it returns exactly $a,z$, causing the
source cache to shrink. The source cache never grows,
and $z$ is retained. Thus, within $n+1$ events, either
a gain occurs or a path vertex is contacted.

When a path vertex is contacted, its reply transfers its
next path head to $a$, strictly shortening the path.
This holds both for $z$ and for a later internal vertex:
its cache contains $a^0$ and that next head. The remaining
suffix is unchanged, since neither endpoint of the exchange
is on the suffix. Off-path contacts had left the original
path intact. Repeat with the shortened path. Its initial
length is at most $n-1$, proving the bound.
\end{proof}

\begin{lemma}[Extremal fresh-copy states]
\label{lem:c2extremal}
In any closed communicating class, fix identity $a$ and
maximize $k=|\{v\ne a:a^0\in C_v\}|$ over its states.
Then $k\ge1$, and every maximizing state has the following
properties:
\begin{enumerate}
\item Every vertex directed-reachable from $a$, other
than $a$ itself, belongs to $I=\{v\ne a:a^0\in C_v\}$.
\item There is no positive-age descriptor for $a$ anywhere.
Thus $I$ is exactly the set of peers storing identity $a$.
\end{enumerate}
Moreover a maximizing state with $C_a=\varnothing$ exists
in the same closed class.
\end{lemma}
\begin{proof}
The integer maximum is attained because its possible
values are $0,\ldots,n-1$. The first property follows from
Lemma~\ref{lem:c2sourcegain} and closedness.

For the second, suppose $p$ stores $a^A$ with $A>0$.
If its oldest partner is $a$, contact $a$. No old zero
holder is touched, and the owner now caches the nonzero
holder $p$. Lemma~\ref{lem:c2sourcegain} increases $k$.
Otherwise $C_p=\{a^A,u^B\}$ with oldest partner $u$.
The request $\{p^0,a^{A+1}\}$ forces the responder cache
to be $\{p^0,a^b\}$. If $u$ previously stored $a^0$,
it keeps that zero and its full reply also gives $p$ a
zero, a contradiction. Otherwise $b>0$ and neither $p$
nor $u$ is a zero holder. Activate $u$: its uniquely
oldest partner is $a$. The owner learns fresh $u^0$,
with no old zero holder affected, and again the source-only
gain lemma contradicts maximality. This proves both properties.
They also imply $k\ge1$: if $k=0$, the absence of positive
copies means no peer caches $a$. Weak connectivity then
forces $C_a$ to be nonempty, and an owner activation
creates a protected $a^0$ copy, a contradiction.

Let $R$ be the vertices reachable from $a$ in a maximizing
state. Its caches contain only identities in $R$, and
every nonowner in $R$ stores $a^0$. If $|R|\ge3$, freeze
the complement and apply Proposition~\ref{prop:anchorcanonical}
within $R$, with $a$ as the final root. All old $R$
nonowners end with $a^0$ and the owner empty; all outside
caches are untouched. If $|R|=1$ the owner is already
empty. If $R=\{a,b\}$, then $C_b=\{a^0\}$ and the
single exchange $a\to b$ empties $a$ without changing
$b$'s zero. These paths preserve $k$ at their endpoints
and stay in the closed class.
\end{proof}

\begin{lemma}[Four-exchange fresh-copy gain]
\label{lem:c2infection}
Suppose $C_a=\varnothing$, $C_p=\{a^0,x^A\}$,
and $x$ does not store $a^0$. Then four legal exchanges
increase the number of $a^0$ holders, preserving all
previous such holders at the endpoint:
\[
 p\to x,\qquad x\to a,\qquad a\to p,\qquad x\to p.
\]
\end{lemma}
\begin{proof}
The first partner $x$ is oldest or tied with $a^0$.
Its received pair forces $C_x=\{p^0,a^1\}$: if an
$a$ descriptor was already there, its age was a positive
integer and minimum merging with one gives one.
The initiator $p$ cannot retain $x$ after this exchange,
since it removes $x$ and a responder never returns itself.
Next $x$ necessarily contacts $a$, giving
\[
 C_a=\{x^0,p^1\},\qquad C_x=\{p^1\}.
\]
The next owner activation necessarily contacts $p$,
sending $a^0,x^1$, so $C_p=\{a^0,x^1\}$ exactly.
Finally singleton $x$ contacts $p$ and receives $a^0$
and its own identity. It retains only $a^0$; at $p$,
incoming fresh $x^0$ refreshes $x^1$ and requires no
eviction. The final caches are
\[
 C_p=\{a^0,x^0\},\qquad C_x=\{a^0\}.
\]
All other old zero holders were untouched.
\end{proof}

\begin{theorem}[Full uniform sampling for $c=\ell=2$]
\label{thm:thesisc2}
For every $n\ge3$, the adopted CYCLON algorithm with
cache capacity two and shuffle length two converges from
every weakly connected initial state with arbitrary finite
integer ages to a unique full-state stationary law $\pi$.
There exist $T=T(n)<\infty$ and $\gamma=\gamma(n)>0$ such that
\[
 \inf_x\widehat P_2^T(x,\{H(V,2)\})\ge\gamma,
 \qquad
 \sup_x\|\widehat P_2^t(x,\cdot)-\pi\|_{\TV}
 \le(1-\gamma)^{\lfloor t/T\rfloor}.
\]
For every observer and $b=1,2$, stationary availability
is positive and its ordered $b$-sample converges to uniform
over $(V\setminus\{i\})_b$ conditional on availability.
No efficient bound on $T$ or $\gamma$ is asserted.
\end{theorem}
\begin{proof}
We first show that every closed communicating class
contains a universal zero-age anchor for a fixed identity
$a$. Fix such a class and the maximum $k$ of
Lemma~\ref{lem:c2extremal}. Among its maximizing states
with $C_a=\varnothing$, let
\[
 I=\{v\ne a:a^0\in C_v\},\qquad
 J=V\setminus(I\cup\{a\}).
\]
Lemma~\ref{lem:c2infection} rules out every cache edge
from $I$ to $J$: it would increase $k$. The extremal
lemma also rules out any descriptor for $a$ in $J$.
Thus each $I$ cache consists of $a^0$ and at most one
residual identity from $I$.

Choose such a state minimizing the number $d$ of weak
components of the residual graph on $I$. This minimum
exists since $1\le d\le k$. Suppose $J$ is nonempty.
Weak connectivity, the empty owner, and absence of
$I\to J$ and $J\to a$ edges imply that a peer
$P\in J$ caches some member of $I$. Inspect an actual
oldest partner of $P$ and keep $P$ frozen for the next
normalization.

By Lemma~\ref{lem:anchorreduction}, residual exchanges
on $I$ simulate capacity-one singleton shuffles, keeping
all anchors at zero and the owner empty. Within each
residual component, the capacity-one theorem makes it
a zero-age rooted star with any prescribed root.
The component partition and $d$ do not change.
In particular the selected root $Q$ has $C_Q=\{a^0\}$;
every other member of its component has
$\{a^0,Q^0\}$. None of $P$'s ages or oldest choices changes.

If $P$'s inspected oldest partner $u$ belongs to $J$,
its other entry is a target $Q\in I$. Normalize $Q$'s
component with root $Q$, and exchange $P\to u$ using
$Q$ as extra. Protection forces
$C_u=\{P^0,Q^b\}$ for some $b\ge0$; choose $u\to Q$
next, which is oldest or tied. Root $Q$ replies only
$a^0$, leaving $C_u=\{a^0,P^1\}$ while $Q$ loses its
zero. Peer $P$ still has no anchor: both its original
endpoints were in $J$ and their pre-exchange caches had
none. The zero count stays $k$, the owner stays empty,
and new holder $u$ has an edge to nonholder $P$.
Lemma~\ref{lem:c2infection} contradicts maximality.

Otherwise $P$'s inspected oldest partner is $Q\in I$.
Normalize its component with root $Q$ and exchange
$P\to Q$. If $P$ is a singleton, it receives $a^0$
while the one incoming fresh $P$ fits alongside the
root's $a^0$, increasing $k$ directly.
Otherwise write $C_P=\{Q,X\}$, with $X\ne a,Q$.
Root $Q$ replies only $a^0$ and is forced to retain
fresh $P$ and the extra $X$, losing its anchor. The new
zero set is exactly $I'=(I\setminus\{Q\})\cup\{P\}$,
and $C_P'=\{a^0,X^{A+1}\}$ for the old age $A$ of $X$.
If $X\in J$, the fresh-copy gain lemma applies immediately.
If $X\in I$ and $Q$'s former residual component had
another member $v$, that member still has
$C_v=\{a^0,Q^0\}$ with $Q$ now a nonholder, giving
the same contradiction.

The only remaining case is that $Q$ was a singleton
residual component and $X$ belongs to another component.
No other $I$ cache points to $Q$. Removing this isolated
vertex and adding new holder $P$ with residual edge
$P\to X$ reduces the component count to $d-1$.
All residual identities still lie in $I'$, the owner
is empty, and the zero count is still maximal. This
contradicts minimality of $d$. Hence $J$ is empty and
$k=n-1$.

Proposition~\ref{prop:anchorcanonical} now sends every
closed class to the same prescribed labelled atom
$h=H(V,2)$. Therefore there can be only one closed
class. The uniform finite-family reset of
Theorem~\ref{thm:thesisgeneralreset} supplies $F_B$, its
horizon $H$, and probability $\delta$.
By the finite-class argument in
Lemma~\ref{cor:thesisgeneralregimes}, every member
of $F_B$ has a finite positive path to a closed class,
and hence to $h$. Choose one such path for each member.
Their lengths have a finite maximum $L$.
Pad them to length $L$ at $h$'s empty-root self-loop and
let $\varepsilon>0$ be the minimum padded-path probability.
Then
\[
 \widehat P_2^{H+L}(x,\{h\})\ge\delta\varepsilon
 \qquad\text{for every }x.
\]
Lemma~\ref{lem:atom} gives the theorem with
$T=H+L$ and $\gamma=\delta\varepsilon$.
Relabelling equivariance gives the same minorization
for every relabelled canonical atom. For any observer,
one such atom gives it two cached identities, proving
positive stationary availability for $b=1,2$.
Theorem~\ref{thm:stationary} proves conditional uniformity.
\end{proof}

\begin{proposition}[Refresh without reducing the endpoint holder count]
\label{prop:zeroify}
For $\ell=2$ and arbitrary $c\ge2$, any state can reach
a state in which all existing descriptors for a chosen
identity $a$ have age zero, with endpoint holder count
at least its initial $N_a$. Every initially zero copy
is preserved throughout the construction.
If a positive-age copy exists, one block of at most
$T_{n-1}+n+c$ events strictly increases $Z_a$ and has
endpoint $N_a$ at least its starting value. At most
$n-1$ such blocks suffice.
\end{proposition}
\begin{proof}
Freeze owner $a$ and all its zero-age holders, and apply
Theorem~\ref{thm:firstcontact} to a positive copy outside
that frozen set. The controlled path activates only
positive holders and always sends $a$ as extra before
contacting the owner. If it first contacts a zero holder,
that responder's zero is protected and returned to the
initiator, increasing $Z_a$ without reducing $N_a$.

If it contacts owner $a$, the initiating peer $p$ loses
its positive copy, and the owner retains fresh $p^0$.
Only this terminal step can decrease $N_a$, by one.
Now activate only $a$, retaining target $p$ as extra
until it is selected. Whenever the responder holds $a$,
require it to reply with $a$ and with $p$ if known.
The returned self identity leaves at most one usable
reply, so $p$ cannot be evicted and the owner cache
never grows. Each first responder learns and retains
$p$; repeats return $a,p$ and shrink the cache.
Within $n+c$ events some nonholder, at latest $p$, is
therefore contacted. It receives protected fresh $a^0$,
restoring the lost holder and increasing $Z_a$.
Earlier positive responders may also have become zero;
all old zeros were passive and protected.

The first-contact horizon is at most $T_{n-1}$, since
the polynomial $T_s$ increases with $s$. Repeat while
a positive copy remains. Every block increases the
bounded integer $Z_a$, so at most $n-1$ blocks are
possible, and endpoint holder counts never decrease.
\end{proof}

\begin{corollary}[Maximum zero count also bounds all holders]
\label{cor:maxholders}
Fix any closed communicating class at $\ell=2$ and
identity $a$, and let $k=\max_x Z_a(x)$ in that class.
Then $N_a(x)\le k$ for every state of the class.
Every maximizing state has exactly $k\ge1$ holders,
all at age zero.
\end{corollary}
\begin{proof}
Apply Proposition~\ref{prop:zeroify} to any class state.
Its endpoint remains in the closed class and has
$Z_a=N_a\ge N_a(x)$, proving the bound.
At a maximizing state, $k=Z_a\le N_a\le k$ excludes
positive copies. If $k=0$, no peer knows $a$; weak
connectivity forces its owner cache to be nonempty,
and an owner activation creates $a^0$, a contradiction.
\end{proof}

\begin{lemma}[Owner service and empty-owner extremal states]
\label{lem:generalownerservice}
At $\ell=2$ and arbitrary $c\ge2$, if owner $a$ caches
a nonholder of identity $a$, a finite owner-only path
increases $N_a$ without losing any existing holder.
The same conclusion holds if a nonholder is directed-reachable
from $a$. In every closed class there is a state maximizing
$Z_a=k$ with $C_a=\varnothing$.
\end{lemma}
\begin{proof}
For a cached nonholder target $z$, activate only $a$ and
send $z$ as extra until the first nonholder contact.
At a holder responder require $a$ in the reply, and $z$
when known. The self descriptor $a$ is discarded at its
owner, leaving at most one usable reply. Thus the target
cannot be evicted and the owner cache cannot grow.
First visits establish $z$ at responders; repeats return
$a,z$ and shrink the owner cache. Within $n+c$ events
some nonholder, at latest $z$, is contacted. Every old
holder was passive and retained received fresh $a^0$;
the nonholder now accepts a new protected copy.

For a directed path to a nonholder, choose a shortest
one, whose internal vertices are holders. Retain its
first head $z$ as extra during off-path contacts.
Before a gain, the same repeat-and-shrink argument forces
a contact with a path vertex. Its reply can contain
$a$ and the next path head, shortening the path. The
remaining suffix is untouched. Repetition reaches a
nonholder in finitely many events. This extends the
argument of Lemma~\ref{lem:c2sourcegain} to arbitrary
capacity and arbitrary ages of the held $a$ descriptors.

Now take a state maximizing $Z_a=k$ in a closed class.
Corollary~\ref{cor:maxholders} says that all its holders
are zero holders, and $N_a\le k$ throughout the class.
The preceding path argument implies that the set $R$
reachable from $a$ is contained in $I\cup\{a\}$, where
$I$ is the zero-holder set. All caches in $R$ contain
only identities in $R$. Freeze its complement. For
$|R|\ge3$, Proposition~\ref{prop:anchorcanonical}, with
effective capacity $\min(c,|R|-1)$ and root $a$, leaves
the owner empty and all other $R$ peers with $a^0$.
The effective capacity changes no transition: caches
have no identities outside $R$, so a larger configured
capacity cannot cause an additional overflow.
If $R=\{a,b\}$, one exchange $a\to b$ suffices,
since $C_b=\{a^0\}$; if $R=\{a\}$ nothing is needed.
Outside zero holders are untouched. The endpoint is
therefore an empty-owner maximizing state in the same class.
\end{proof}

\begin{lemma}[Outward holder gain when $n\le2c$]
\label{lem:halfdensegain}
Suppose $\ell=2$, $c\ge3$, $n\le2c$, $C_a=\varnothing$,
and exactly $h$ peers hold identity $a$, all at age zero.
If any holder caches a nonholder, a finite legal path
increases the total holder count beyond $h$.
No connectivity assumption is needed for this lemma.
\end{lemma}
\begin{proof}
We describe a controlled path and stop as soon as the
holder count exceeds $h$. Initially track the specified
holder $p$, and let $x$ be a cached nonholder target.
Activate only $p$ while it contacts other holders.
Choose a nonanchor oldest entry in a tie, send $a$ as
extra, and require $a^0$ and $p$ when known in the reply.
The anchor is refreshed to zero; at most one other usable
identity arrives, so the active cache cannot grow.
Every other resident, including target $x$, is unsent
and cannot be a victim. First visits establish $p$ at
responders; repeats return $a,p$ and shrink the cache.
Consequently a nonholder is contacted within $n+c$ events.
All old holders have been preserved before this contact.

Consider any such contact $p\to q$, with $q$ a nonholder
and outgoing packet $\{p^0,a^A\}$, $A>0$. If $p$ is
underfull, its mandatory deletion creates at least two
vacancies; every reply fits and the old anchor survives,
while $q$ accepts a new one. If $C_q$ contains an identity
in $\operatorname{ids}(C_p)\cup\{p\}$, require one such
identity in the reply. It is duplicate or self-filtered
at $p$ (and cannot be $q$), leaving at most one usable
new identity. Again the old anchor survives and a gain
occurs. These reply choices have legal victims: when
fresh $p$ is known at $q$ only $a$ is new there; otherwise
both incoming identities are new and all returned
residents are eligible. The prescribed shared reply
does not restrict feasibility.

Thus the only remaining case has $p$ full, two usable
novel reply identities, and
\[
 \operatorname{ids}(C_q)\cap
       (\operatorname{ids}(C_p)\cup\{p\})=\varnothing.
\]
In particular $|C_q|\le n-c-1\le c-1$.
The active cache must evict its only extra, $a$, and
remains full; the passive cache accepts protected $a$
and fresh $p^0$, ending with size $\min(c,|C_q|+2)$.
Call this event a migration and track the new holder $q$.
The old holder $p$ is now a nonholder, cached as a fresh
predecessor at $q$. The total holder count is still $h$;
the only possible positive-age $a$ descriptor is at the
tracked holder. All other holders remain at zero.

After a migration, take one activation of the tracked
holder $u$, keeping all other initiators frozen. If its
partner is a zero holder, send $a$ and require $a^0$ in
the reply, with $u$ when known. Its cached nonholder
predecessor is unsent and survives. The anchor becomes
zero, the cache cannot grow, and the preceding
single-holder service applies. If the partner is a
nonholder, use the same gain-or-migration choice above.

The remaining possibility is partner $a$ itself. Let $v$
be the cached nonholder predecessor and send it as extra.
Because the owner has stayed empty, this gives
\[
 C_a=\{u^0,v^{B+1}\},\qquad B\ge0,
\]
and $u$ loses its anchor. The next owner activation must
contact $v$; send $u$ as extra. Peer $v$ acquires fresh
$a^0$, restoring the holder count to $h$. At the owner,
its one remaining extra $u$ and at most two reply entries
fit in capacity $c\ge3$, so $u$ is retained. Peer $u$
is still a nonholder. Lemma~\ref{lem:generalownerservice}
now supplies one more holder. This double callback
therefore gives a strict gain even though the first
exchange temporarily lost one holder.

If the tracked cache is ever underfull, it must lead
to a gain: its next contact either uses this owner
callback, gains at a nonholder without extra eviction,
or refreshes the anchor at a zero holder. In the last
case its cache remains underfull, and single-holder
service reaches a nonholder while retaining the anchor.

It remains to rule out infinitely many full-cache
migrations. Until a gain or a transition to an underfull
tracked cache, count full caches at all nonowners.
Passive cache cardinalities never decrease. An active
full tracked cache either stays full or invokes the
underfull finish just described. At a migration, its
nonholder responder was underfull; if the new tracked
cache is full, this strictly increases the count.
This integer is at most $n-1$, so only finitely many
such migrations can occur. Between migrations,
single-holder service takes finitely many events.
The path must therefore finish with a strict holder gain.
Every prescribed choice is an existing positive-probability
outcome of the adopted protocol.
\end{proof}

\begin{lemma}[Outward holder gain for every capacity $c\ge3$]
\label{lem:generaloutwardgain}
Suppose $\ell=2$, $c\ge3$, $C_a=\varnothing$, and
exactly $h$ peers hold identity $a$, all at age zero.
If a holder caches a nonholder, a finite legal path
increases the total holder count beyond $h$.
No restriction on $n/c$ or initial ages is needed.
The path length in this lemma may depend on the ages.
\end{lemma}
\begin{proof}
Use the controlled trajectory in the proof of
Lemma~\ref{lem:halfdensegain}, replacing its full-cache
potential by the argument below. The terminal gadgets
in that proof did not use $n\le2c$: an underfull
current holder with a cached nonholder leads to a gain;
contact with the empty owner gives a gain by double
callback; and a nonholder responder sharing any identity
with the current cache or current self gives a gain
by a suitable reply. We may therefore restrict attention
to a full current holder, with its nonholder predecessor
or initial target still cached, until one of these
terminal cases occurs.

Only the current holder initiates. At every nonowner
contact it sends its fresh self and $a$. At a nonholder
contact not giving a gain, the responder cache is
disjoint from the current cache and current self.
Both reply identities are novel; the active holder
loses its anchor and stays full, and the responder
becomes the new current holder. At a zero-holder contact,
the reply contains $a^0$, so the active holder remains
current and its anchor is refreshed to zero.

Let $Z$ be the initial zero holders other than the
initial current. Throughout an unfinished trajectory,
the holder set is exactly $Z\cup\{\text{current}\}$.
Every member of $Z$ is permanently passive and keeps
its zero. Every past current is a nonholder until it
becomes current again in a migration. The cached
nonholder predecessor is unsent at zero-holder contacts
and therefore survives them. These facts justify the
terminal gadgets at every later state as well.

Initially mark every descriptor old. Mark fresh self
descriptors young and use young-OR on duplicate merges,
independently of which input supplies the minimum
numerical age. Let $F$ count young cache entries,
so $0\le F\le nc$. Identity $a$ stays old because
its owner is frozen. Every identity in $Z$ also stays
old: none of these peers ever initiates, so no fresh
self descriptor with that identity is generated.

We can make $F$ nondecreasing, with a strict increase
at each zero-holder contact. Suppose current $p$
contacts $q\in Z$. If $q$ knows $p$, reply $a,p$;
the active cache loses its mandatory partner and gains
no new identity, becoming underfull and invoking the
terminal gain gadget. Otherwise fresh $p$ is new at $q$.
Require reply $a,D$ if another resident exists.
If $q$ is underfull no victim is needed, and any such
$D$ is allowed. If $q$ is full, choose a nonanchor
$D$ absent from $C_p$ and use it as the responder's
single victim. Such a $D$ exists: its $c-1$ nonanchor
residents can share at most $c-2$ identities with $C_p$,
which also contains $a$ and the selected identity $q$.
Neither of those two identities is a nonanchor resident
of $q$. The reply is legal, protects the active anchor,
and copies a young victim to a previously absent
location at $p$ if compensation is needed. The selected
partner $q$ is old, and no young initiator extra can
be evicted since its only extra is old $a$.
Fresh $p$ supplies an additional young entry at $q$.
Thus $F$ strictly increases.

At a nonholder migration, choose every required
responder victim inside the two-entry reply. The
disjointness means that every young victim is copied
to a new location at the initiator. Its only extra
victim is old $a$, and fresh $p$ adds a young entry
at the responder. Only the selected partner can
remove a young initiator resident. Hence
\[
 F'-F\ge1-\mathbf1_{\{\text{selected partner is young}\}}.
\]
So $F$ cannot decrease, and it increases at every
migration selecting an old partner. No assumption
that the responder was underfull has been used.

Suppose the controlled path never reached a gain.
Every finite unfinished state admits a next event
of the prescribed kind. Since $F$ is a bounded
nondecreasing integer, it would eventually be constant.
There are then no zero-holder contacts, and all events
are nonholder migrations selecting young partners.
Start after one such migration, with current anchor
age $A\ge1$, and count subsequent events by $t$.
At time $t$ its age is exactly $A+t$: a migration
increments and transfers it to a previous nonholder,
with no duplicate refresh.

For every identity $s$, let $M_s(t)$ be its largest
stored age anywhere, or zero if absent, and put
$K_s(t)=M_s(t)-t$. Each event increments only active
residents by one; fresh selves have age zero, and
copies and minimum merges cannot exceed their input
ages. Thus every $K_s$ is nonincreasing. More strongly,
when $s\ne a$ is selected as partner, its sole active
copy is removed, it is not sent as extra, and its own
responder cannot return itself. All remaining $s$
copies have unchanged or reduced ages. Therefore
\[
 K_s(t+1)\le K_s(t)-1
 \quad\text{whenever partner }s\text{ is selected}.
\]
The oldest rule simultaneously requires the selected
descriptor's age to be at least $A+t$, giving
$K_s(t)\ge A$. Each identity can consequently be
selected at most $(M_s(0)-A+1)_+$ times during this
constant-$F$ period. Summing over the finitely many
$s\ne a$ bounds its length, contradicting an infinite
unfinished path. Hence some terminal gain occurs.
Every prescribed exchange has positive probability,
so its finite controlled path is accessible.
\end{proof}

\begin{theorem}[Full uniform sampling for every two-entry shuffle]
\label{thm:thesisgeneraltwo}
For every $n\ge3$ and $2\le c\le n-1$, the adopted
algorithm with fixed shuffle length two converges
from every weakly connected initial state with
arbitrary finite integer ages to a unique stationary
full-state law $\pi$. There exist finite $T(n,c)$
and positive $\gamma(n,c)$ such that
\[
 \inf_x\widehat P_2^T(x,\{H(V,c)\})\ge\gamma,
 \qquad
 \sup_x\|\widehat P_2^t(x,\cdot)-\pi\|_{\TV}
       \le(1-\gamma)^{\lfloor t/T\rfloor}.
\]
For every observer and $1\le b\le c$, stationary
availability is positive and its ordered local
$b$-sample converges to uniform conditional on availability.
No practical mixing-time bound is asserted.
\end{theorem}
\begin{proof}
Capacity two is Theorem~\ref{thm:thesisc2}.
For $c\ge3$, argue as follows.
Fix a closed class and identity $a$. By
Corollary~\ref{cor:maxholders} and
Lemma~\ref{lem:generalownerservice}, it contains an
empty-owner state maximizing $Z_a=k$, with exactly $k$
zero holders and no positive copies. Put
$I=\{v\ne a:a^0\in C_v\}$ and
$J=V\setminus(I\cup\{a\})$. No class state can have
more than $k$ total holders. Lemma~\ref{lem:generaloutwardgain}
therefore rules out every edge from $I$ to $J$.
All residual identities in $I$ are consequently in $I$.

If $J$ is nonempty, weak connectivity, the empty owner,
and absence of $J\to a$ and $I\to J$ edges give an
edge from $J$ to $I$. Freeze $I\cup\{a\}$ and use
Theorem~\ref{thm:firstcontact} to reach, but pause just
before, an outside initiator $P$ contacts $Q\in I$.
No outside cache contains $a$ before this first contact.
Pause all outside peers. Reserve the zero anchor at
every member of $I$ and use
Lemma~\ref{lem:anchorreduction} to simulate a
singleton clearing block at $Q$, with residual capacity
$c-1$. All its residual partners lie in $I$.
Stop when the residual cache empties, before it contacts
the owner. This leaves $C_Q=\{a^0\}$, preserves every
other zero holder, and leaves $P$'s cache and actual
oldest choice of $Q$ unchanged.

Resume $P\to Q$. Its at-most-two request identities
and $Q$'s anchor fit in capacity $c\ge3$, so $Q$
keeps its zero. Its sole reply gives $P$ another zero,
contradicting maximality of $k$. Hence $J$ is empty:
every closed class contains a universal zero anchor.

Proposition~\ref{prop:anchorcanonical} sends every closed class to
the same prescribed labelled $H(V,c)$.

Although the new outward path can have an
age-dependent length, the finite-family argument
still gives a uniform minorization: the finite reset
family $F_B$ has only finitely many chosen paths to
this common atom. Their lengths have a finite maximum. Pad at the
empty-root self-loop, take the positive minimum probability
of the padded paths, and apply the uniform reset,
as in Theorem~\ref{thm:thesisc2}. The same argument
gives uniqueness and geometric convergence.
Relabelling symmetry yields conditional uniformity;
relabelled canonical atoms give positive availability
for every $b\le c$.
\end{proof}

\section{Three-entry shuffles: the third induction base}

\begin{theorem}[Uniform first contact for three-entry shuffles]
\label{thm:threefirstcontact}
At $\ell=3$, fix a nonempty proper
$G\subset V$, put $S=V\setminus G$, $s=|S|$, and
suppose some outside cache contains an identity $r\in G$.
There are constants $T_3(n,c,s)<\infty$ and
$\eta_3(n,c,s)>0$, independent of all initial ages,
such that an $S$-initiator path contacts $G$ within $T_3$
events. Under the original scheduler the first exchange
across the cut has the geometric bound
\[
 \Pr_x\{\tau>jT_3\}\le(1-\eta_3)^j.
\]
The constructed path always activates an $r$ holder,
sends $r$ as extra when the partner is not its owner,
and requires $r$ in the reply whenever known. Its total
$r$-holder count cannot decrease except at a terminal
contact with owner $r$.
\end{theorem}
\begin{proof}
Keep $G$ inactive until contact. Let $m$ be the minimum
outside age of $r$ and $A\subset S$ its minimum-age holders.
For an exchange internal to $A$, reserve $r$ as one
request extra and one reply entry. Removing this anchor
gives the exact length-two kernel with capacity $c-1$:
the minimum rule restores the anchor to $m$, its incoming
copy protects it, and all other updates are unchanged.
There is only an oldest-partner qualification. If the
residual oldest age is at least $m$, a reduced oldest
partner is an actual oldest partner. If all residual
ages are smaller than $m$, or the residual cache is
empty, the actual next partner is owner $r\in G$,
which is immediate success. No ages are shifted or truncated.

Suppose a residual $A$ cache has an identity outside $A$
other than $r$. On $V\setminus\{r\}$ use
Theorem~\ref{thm:firstcontact} with active set $A$;
put arbitrary empty auxiliary caches at its complement.
Until the terminal contact, only $A$ caches are read,
so those auxiliary choices are immaterial. Simulate
this lower-length path until a contact with $G$, or
until an initiator in $A$ contacts $q\in S\setminus A$.
In the latter event choose the actual reply to contain
$r$ whenever known. Its age is greater than $m$, while
the sent anchor has age $m+1$. At least one minimum
copy disappears, and no new minimum copy is created.
The auxiliary terminal reply is never needed.

Otherwise all residual $A$ identities lie in $A$.
Choose one weak component $K$ of the residual graph.
For $|K|=1$ its cache is just the anchor and contacts
$G$ on activation. For $k=|K|\ge2$ use effective
capacity $d_k=\min(c-1,k-1)$ and length
$e_k=\min(2,k-1)$. Its lower-length kernel has a
uniform common-atom horizon $U_k$. Indeed, $k=2$
uses length one, and all other capacities are covered
by Theorem~\ref{thm:thesisgeneraltwo}.
Simulate a path to a reduced canonical
state with an empty root cache. A reduced empty-cache
activation instead contacts $G$, so any such interruption
is already success. If no interruption occurs, activate
the empty reduced root once to contact $G$.

Let $A_b$ be the length-two first-contact horizon on
$n-1$ peers at capacity $c-1$ with active-set size $b$,
and choose
\[
 E=1+\max\bigl(\{A_b:1\le b\le n-2\}
        \cup\{U_k:2\le k\le n-1\}\cup\{0\}\bigr).
\]
Within $E$ events one minimum copy disappears or $G$
is contacted. All preterminal initiators and responders
in this phase store $r^m$, so their other minimum copies
are retained unchanged. At most $s$ such phases raise
the minimum outside age by at least one, or contact $G$.

Put $R_s=(sc+1)c$. Raise the outside minimum $r$ age
$R_s+1$ times using the preceding phases, unless $G$
is contacted earlier. Every outside $r$ age then
exceeds $R_s$. Mark all current descriptors old and
future fresh selves young, with young-OR duplicate
colors. For the next $R_s$ events every young age is
at most the elapsed number of events. Thus an oldest
partner of an $r$ holder is old throughout this horizon.

At each step choose an $r$ holder $p$ maximizing its
young count $f_p$ among outside $r$ holders. Use the
reply/victim choices of Lemma~\ref{lem:generalcolor},
but reserve old $r$ among its extras and require all
$Q\cap U$ in the reply, in addition to the responder
victims $D_q$. Reserving $r$ leaves room for the
required young difference extra at length three:
if the cache has only two entries, both its anchor
and selected partner are old and there is no young
difference to send. Requiring all $Q\cap U$ is feasible
by the same bound $|D_q|+|Q\cap U|\le\min(3,|Q|)$.
In particular an existing responder $r$ is returned.
The young-count estimate in that lemma still gives
$F'\ge F$, where $F$ counts young entries in $S$.

Write $M=\max\{f_v:v\in S,\ r\in C_v\}$.
If $F'=F$, the same calculation gives
$Y_p\subseteq Y_q$, $p\in Y_q$, and $f_q\ge f_p+1$.
Thus $q$ could not already have held $r$, by maximality
of $f_p=M$. Its new $r$ copy is protected. On this
plateau no young initiator copy is lost; only $t$
young responder victims transfer to new locations
at $p$, with
\[
 f_p'=f_p+t,\qquad f_q'=f_q-t,\qquad
 0\le t\le f_q-f_p-1.
\]
Consequently the new $r$ holder $q$ has
$f_q'\ge M+1$, even if $p$ loses its anchor.
The pair $(sc-F,c-1-M)$ strictly decreases
lexicographically at every nonterminal step.
There are only $(sc+1)c=R_s$ possible pairs,
since every $r$ holder contains at least one old
descriptor. Therefore $G$ is contacted within $R_s$
such steps.

The minimum-raising and final routing phases activate
only $r$ holders and send $r$ whenever the partner is
not its owner. Existing responder copies are returned
and protected; a newly created responder copy compensates
any initiator loss. This proves the holder-count claim.
One may take
\[
 T_3=sE(R_s+1)+R_s,\qquad
 \eta_3=(nc\,16^c)^{-T_3}.
\]
The probability estimate is the all-length controlled
choice bound. Before the first crossing, $G$-internal
events cannot change the outside caches, and outside
internal events preserve some $r$ there. Restarting
the one-block bound proves the geometric tail.
\end{proof}

\begin{proposition}[A universal zero anchor suffices at length three]
\label{prop:threezeroanchor}
At $\ell=3$ and $3\le c\le n-1$, a state in which every
nonowner stores $a^0$ has a finite path to any prescribed
canonical state $H(V,c)$. Hence, if every closed class
can reach a universal zero anchor, full convergence and
conditional uniform sampling follow at length three.
\end{proposition}
\begin{proof}
Freeze $a$ and reserve its zero anchor in every request
and reply. Lemma~\ref{lem:anchorreduction} simulates
the length-two, capacity-$c-1$ kernel on residual
components of $V\setminus\{a\}$; reduced empty-cache
activations are omitted. The lower-length common-atom
theorems apply to every component by exactly the
parameter check in Theorem~\ref{thm:threefirstcontact}.

To empty the owner, inspect an actual oldest $q\in C_a$
and keep $a$ frozen while its residual component is
normalized with root $q$. This gives $C_q=\{a^0\}$,
without changing $a$'s oldest choices. Activate $a\to q$.
Its sole reply is the self descriptor $a$, so $C_a$
loses $q$ and no other entry. At $q$, incoming fresh
$a^0$ protects its anchor, and at most two extras fit
beside it in capacity $c\ge3$. All anchors remain zero.
At most $c$ repetitions empty the owner.

Normalize each residual component. If there is more
than one, choose two roots $u,v$, each with the singleton
cache $\{a^0\}$. The same four exchanges used in
Proposition~\ref{prop:anchorcanonical},
\[
 u\to a,\quad v\to a,\quad a\to v,\quad a\to u,
\]
end with $C_a=\varnothing$, $C_u=\{a^0\}$ and
$C_v=\{a^0,u^0\}$. All involved pre-exchange caches
have size at most two, so their packets and replies
are identical at lengths two and three. Other anchors
are unchanged. This merges two residual components.
Repeat normalization and merging, then normalize the
connected residual graph once more. This gives the
canonical state with root $a$. The root-change gadget
of Proposition~\ref{prop:anchorcanonical} also uses
only caches of size at most two, so any labelled
canonical order is reachable.

If universal zero anchors are accessible from every
closed class, all classes therefore share this atom.
The finite-reset and symmetry argument of
Theorem~\ref{thm:thesisc2} gives the final implication.
\end{proof}

\begin{lemma}[A designated owner target can be retained]
\label{lem:ownertargetretain}
For any shuffle length $\ell\ge2$, suppose owner $a$ caches a
designated target $p$ and contacts some $q\ne p$.
Choose $p$ as extra and require $a$ in the reply,
assuming the responder has it. There is a legal
initiator-victim choice retaining $p$. This is a
positive-probability controlled choice, rather than
a claim about every uniform victim outcome.
\end{lemma}
\begin{proof}
If $p$ is returned, received-identity protection suffices.
Otherwise let $d=|C_a|$ before partner removal and
$e=\min(\ell-1,d-1)$ be its number of extras.
The returned self descriptor leaves at most $\ell-1$
usable identities. If $h$ of those are resident
duplicates, at most $\ell-1-h$ are new. The number
of required victims is at most
\[
 (\ell+d-c-2-h)_+.
\]
Among extras other than $p$, at least $e-1-h$ are
not protected by incoming identities. Since
\[
 \ell+d-c-2\le e-1
\]
both when $d\ge\ell$ (use $d\le c$) and when
$d<\ell$ (use $\ell\le c$), these other extras
suffice. If the displayed lower bound is negative,
the victim requirement is zero. Thus $p$ need not
be selected as victim.
\end{proof}

\begin{lemma}[Learned-target owner boundary service at capacity three]
\label{lem:c3learnedcallback}
Let $c=\ell=3$. Fix owner $a$, and let $H$ be any subset
of its current zero holders. If $a$ caches an identity
$p\notin H\cup\{a\}$, a finite legal controlled path ends
with an owner contact outside $H$, preserving every original
$a^0$ copy throughout and never decreasing $N_a$.
There is no initial shared-secondary or age-band assumption.
For $H$ equal to the full zero-holder set, the path strictly
increases $Z_a$. No age-uniform numerical horizon is asserted.
\end{lemma}
\begin{proof}
Let only $a$ initiate at first, preferring outside-$H$ heads
in oldest ties, and stop at such a contact. Until then every
responder returns $a^0$ in its complete reply and receives
protected fresh $a^0$. Retain $p$ as an extra. Self filtering
leaves at most two usable reply identities and mandatory
deletion supplies at least one vacancy. If a full owner
requires one victim, its other extra $r$ is eligible: were
$r$ received, at most one usable reply identity would be
novel and no victim would be required. Thus $r$ can be
chosen instead of $p$. An underfull owner needs no victim.
Every contacted responder learns $p$ and keeps it at every
later owner contact. The set of such responders grows only
finitely often; all zero holders stay zero before termination.

Suppose this owner-only path continued forever without a
boundary contact. This is a proof device for finding a finite
path, not an event assigned positive probability. After its
last new learning, each future responder has $a^0,p$ and
at most one residual entry. At most that residual is novel
at the owner, so no further active eviction occurs and owner
size cannot grow. Its number of heads other than $p$ therefore
eventually stabilizes at one or two. Secondary ages in the
visited tail and at the owner have a finite bound $M$:
each exchange replaces endpoint $p$ ages $(u,v)$ by
$\min(u+1,v)$ at both, which lies between the old values
for integer ages.

We first record a local completion step. Suppose a finite
set $S$ consists of zero holders, all knowing $p$, with
every residual identity at $S$ in $S$, and the owner cache
is $\{p,Q\}$ for $Q\in S$. If $p$ is oldest, contact it
and finish. Otherwise $Q$ is an actual oldest head. Freeze
the owner and let only $Q$ initiate, reserving $a,p$, while its residual partner
is in $S$ and its age is at least the current $p$ age.
These are actual oldest choices. First visits install fresh
$Q$ pointers and a repeat self-filters the returned pointer;
thus within $|S|$ exchanges the residual empties or has age
$b<v$, where $v$ is the current $p$ age. All anchors remain
protected. Contact $Q$ from the frozen owner. There is no
active victim and at most one residual owner head remains.
Its secondary age becomes $w=\min(u+1,v)$, where $u$ was
the owner secondary age. If another nonfinishing macro is
needed, its sole residual must satisfy $b>w$; together with
$b<v$ this forces $w=u+1$. Hence every continuing macro
strictly increases the owner secondary age, while its finite
global maximum on $S\cup\{a\}$ cannot increase. Finitely
many macros therefore finish at $p$. All other peers are
frozen in this local construction.

In a stable one-head tail, owner requests contain only
$a,p$, both already known by its responders. Their residual
identities and numerical ages consequently stay fixed. An
empty or outside residual finishes; otherwise the owner's
head eventually follows a directed cycle $S$ of visited
responders. It is residually closed and all its members
know $p$. The local completion step applies.

For a stable two-head tail, a head outside $H$ cannot remain
unselected forever: it would be retained, every visited
responder would acquire it as residual, and a repeat would
return it as a duplicate, removing the last $H$ head. Thus
both heads lie in $H$. Let $S$ be the identities selected
infinitely often. Eventually both heads lie in $S$ as well:
a head outside $S$ appearing after the last such selection
would remain forever, and the same repeat argument would
coalesce the heads. After all members of $S$ are visited
again, their residual pointers also lie in $S$.

Write an update's old owner heads as $q:A,s:B$, with $q$
selected, and its old responder residual as $r:C$. In the
stable case $r\in S$ and $r\ne s$. The new owner heads are
$s:(B+1),r:C$, and the new responder residual is $s:(B+1)$.
In the multiset of both owner-head ages and all residual
ages at $S$, this removes $A,B$, inserts $B+1$ twice, and
merely moves $C$. Its minimum cannot decrease and no copy
at the current minimum is created. Every passive minimum
eventually moves to an owner head, since its source is
selected infinitely often; every owner minimum disappears
on the following update. Thus each minimum vanishes after
finite time, and the restricted integer minimum tends to
infinity.

At a finite time all residuals at $S$ have age at least
the secondary bound $M$. Freeze the owner before its next
actual oldest head $Q\in S$. A $Q$-only residual clearing
path inside $S$ is now legal: each unvisited pointer keeps
its old age at least $M$, every secondary stays at most
$M$, and a repeated responder returns self-$Q$ and clears
the residual. The owner's still-legal contact to $Q$ then
returns only the anchors and deletes one owner head. The
owner has $\{p,R\}$ with $R\in S$. Residual closedness
persists: preparation installed pointers only to $Q$, and
the owner contact installs only $R$ at $Q$. The local
one-head completion step finishes.

Every infinite alternative has supplied a finite prefix
and a finite finishing intervention. All preterminal
responders and preparation endpoints retain their protected
zeros; other original holders are untouched. The final
responder receives protected fresh $a^0$. Every prescribed
finite choice has positive probability under the unchanged
protocol, proving the statement.
\end{proof}

\begin{corollary}[Capacity-three zeroification and nonholder recovery]
\label{cor:c3zeroify}
At $c=\ell=3$, if the owner caches a nonholder of identity
$a$, a finite path increases $N_a$ by at least one and
preserves every existing holder and every existing zero copy. Every state has a finite
path making all its remaining identity-$a$ copies zero,
with endpoint holder count at least its initial $N_a$.
Consequently, in a closed class with $k=\max Z_a$, every
class state satisfies $N_a\le k$, and all $k$ holders at
a maximizing state are zero holders.
\end{corollary}
\begin{proof}
Retain a fixed cached nonholder target $p$. Apply
Lemma~\ref{lem:c3learnedcallback} with the current full zero
set. If the terminal responder is already a positive holder,
its complete reply contains self-$a$, so the same victim
argument retains $p$ at the owner. The zero count has grown,
no holder was lost, and $p$ is still a nonholder. Repetition
must contact some nonholder after finitely many zero gains,
since the fixed $p$ cannot remain excluded from a strictly
growing zero set forever. This is some new holder, not a
claim of service to a prescribed nonholder.

For zeroification, freeze the owner and all zero holders.
From a positive $a$ copy apply the holder-only construction
of Theorem~\ref{thm:threefirstcontact} to the cut containing
the owner and zero holders, using target identity $a$.
A boundary zero-holder contact refreshes the active copy
and protects the responder's zero, increasing $Z_a$ without
decreasing $N_a$. An owner contact loses one active copy but
gives the owner the fresh address of that now-nonholder.
The first assertion restores the lost holder count while
gaining a zero and preserving all previous zeros. Each
stage strictly increases $Z_a$, so finitely many stages
leave no positive copy. Applying this inside a closed
class proves $N_a\le k$; equality of $Z_a$ and $N_a$ at a
maximizer gives the last assertion.
\end{proof}

\begin{lemma}[Owner boundary service at every capacity for length three]
\label{lem:threegeneralcallback}
Let $\ell=3$ and $c\ge3$. Fix owner $a$, a subset $H$ of
its zero holders, and a cached target $p\notin H\cup\{a\}$.
A finite legal path ends with an owner contact outside $H$,
preserving every existing holder of identity $a$ and every
existing zero copy. All preterminal exchanges are internal to
$H\cup\{a\}$, and $p$ stays cached at the owner before the
final contact. Its terminal reply may be chosen freely subject
to the model. No initial shared-target or age-band assumption
is needed. No uniform numerical horizon is asserted.
\end{lemma}
\begin{proof}
Until termination, let only $a$ initiate and prefer outside-$H$
partners in oldest ties. Include $p$ in every request. At
preterminal contacts require $a$ in replies; Lemma~\ref{lem:ownertargetretain} retains $p$
at the owner, even before its responder knows $p$. Every
visited responder learns $p$ and thereafter keeps it. Once
$p$ is known, choose reply $\{a,p,C\}$ with one residual $C$
if present. The owner then receives at most one novel identity,
so its head count $k$ excluding $p$ cannot increase.
If the responder has no residual, or has a residual already
among the owner's retained heads, return that residual when
present: $k$ decreases. Otherwise return a minimum-age residual.
A first learning may increase $k$ by at most one.

For $k\ge2$, send $p$ and another retained head $B$ as extras.
Use a bounded-fair choice of $B$: assign new heads their arrival
time, choose the retained head with oldest last-sent timestamp,
and update that timestamp. Thus any head that stays cached and
is never selected as partner is eventually sent. This is a
controlled policy, not a fairness assertion about the scheduler.

Suppose an owner-only tail never terminates, has no further
learning, and has stable $k\ge2$. At every late event the
responder's residual identities are disjoint from the retained
owner heads, since otherwise the overlap-priority reply would
decrease $k$. In particular $B$ is novel at the responder.
An underfull responder grows by one; a full responder replaces
its returned residual $C$ by $B$, the only eligible residual
victim. Its other identities remain. These passive occupancies
never decrease, so only finitely many growth events occur.
Every subsequently selected responder is therefore full.

Let $S$ be the identities selected infinitely often. Eventually
all owner heads lie in $S$. Indeed, a head outside $S$ appearing
after the last outside-$S$ selection would remain forever.
The extra-selection policy sends it to some recurrent responder,
which retains it until its next visit. That visit has an overlap
and decreases $k$, a contradiction. After this finite prefix,
consider the multiset of the $k$ owner-head ages and all residual
ages at $S$. An update removes the selected root age, increments
all retained root ages, moves the returned residual age to the
root, and inserts a second copy of the chosen extra's incremented
age at the responder. No new copy of the current minimum is
created. Every root copy of that minimum disappears next step;
every passive copy eventually moves to the root, because its
source is recurrent and a minimum residual is returned. Thus
all copies of each integer minimum vanish after finite time,
and the multiset minimum tends to infinity.

A residual at $S$ pointing outside $S$ keeps its fixed age until
returned to the owner. Divergence forces its return, contradicting
eventual owner-head containment in $S$. Hence on a sufficiently
late tail $S$ is residually closed. All its members know $p$.
The $p$ ages at $S$ and the owner have a finite bound $M$:
each endpoint pair is updated to $\min(u+1,v)$, no greater than
the old maximum. Freeze the owner once every $S$ residual has
age at least $M$, before its next actual oldest head $Q\in S$.

Let only $Q$ initiate, reserving $a,p$ and choosing an oldest
residual partner. Require reply $a,p,Q$ when $Q$ is known,
and otherwise $a,p$ plus a residual if present. Each first
visit installs fresh $Q$. A repeat returns self-$Q$, reducing
$Q$'s residual head count; first visits cannot increase that
count. Initial unvisited residuals have age at least $M$,
retained active residuals only increase, and $p$ stays at most
$M$, so every chosen residual is an actual oldest choice.
Residual closedness keeps all contacts in $S$. In finitely many
first visits and head deletions the residual cache empties.
The frozen owner's still-legal contact to $Q$ returns only
$a,p$ and decreases $k$. All old zeros and all learned $p$
entries remain protected.

For a one-head owner phase, the cache is $\{p,Q\}$. If $p$
is oldest, finish. Otherwise $Q$ is oldest. Owner requests
contain only $a,p$, both duplicates at learned responders,
so their residual graphs and numerical ages stay fixed.
Return residuals in a fair cyclic rotation at each responder.
On an infinite stable learned tail, every outgoing edge of an
infinitely visited vertex is then traversed infinitely often.
Its recurrent set $S$ is residually closed and all members know
$p$; an empty residual would instead finish immediately.

Freeze the owner and perform $Q$-only residual clearing as above,
but stop when its residual cache empties or its maximum residual
age is below its $p$ age. The same finite first-visit/repeat
argument applies; every executed choice is an actual oldest
residual. Contact $Q$ from the still-frozen owner. If a residual
of age $b$ is returned, preparation stopped with $b<v$, where
$v$ is the prepared $p$ age. The owner secondary becomes
$w=\min(u+1,v)$. A further nonfinishing macro requires $b>w$,
forcing $w=u+1$. Each continuing macro therefore strictly
increases the owner secondary, bounded by the finite global
$p$ maximum on $S\cup\{a\}$. Only finitely many are possible.

These arguments provide finite prefixes and finite interventions,
not positive-probability infinite schedules. Across restarts the
learned set is monotone, including during clearing. All possible
increases of $k$ are charged to first learnings, at most one per
new learner. Therefore only finitely many clearing-induced head
decreases can occur before a one-head completion or boundary
contact. All preparation endpoints protect their $a^0$ copies;
other existing holders are untouched. The final responder receives
protected fresh $a^0$. Its final reply can be any legal subset;
no earlier argument constrains that terminal choice. Every finite
controlled choice has positive probability, proving the claim.
\end{proof}

\begin{corollary}[Zeroification at every capacity for length three]
\label{cor:threegeneralzeroify}
Let $\ell=3$ and $c\ge3$. A cached nonholder at owner $a$
gives a finite path increasing $N_a$ while preserving all
existing holders and zeros. Every state has a finite path
making all identity-$a$ copies zero with endpoint holder count
at least its initial value. In any closed class, if $k=\max Z_a$,
then $N_a\le k$ throughout the class, and all holders at a
zero-maximizing state are zero holders.
\end{corollary}
\begin{proof}
Repeat the proof of Corollary~\ref{cor:c3zeroify}, replacing
its callback by Lemma~\ref{lem:threegeneralcallback}. A terminal
positive-holder reply is required to include $a$, allowing
retention of the fixed nonholder target by
Lemma~\ref{lem:ownertargetretain}. Each such contact increases
$Z_a$, so some nonholder must be reached after finitely many
stages. The holder-only first-contact construction of
Theorem~\ref{thm:threefirstcontact} then zeroifies positive
copies exactly as before. This corollary does not establish
outward holder gain or full convergence at sparse $c\ge4$.
\end{proof}

\begin{lemma}[A full-shuffle triangle gains a holder]
\label{lem:fullshuffletrianglegain}
Let $\ell=c\ge2$, let $C_a=\varnothing$, and suppose $p$
caches $a$ as an oldest entry and a nonholder $r$ as an
oldest entry among its other identities. The three exchanges
$p\to a$, $a\to r$, $p\to r$ are legal and end with
exactly the original identity-$a$ holders together with $r$.
Both $p$ and $r$ end with $a^0$; all other nonowner caches
are unchanged. Other initial holders need not have age zero.
\end{lemma}
\begin{proof}
Write $C_p=\{a^A\}\cup\{v^{b_v}:v\in E\}$, where
$A\ge\max b_v$ and $b_r=\max b_v$. After $p\to a$,
\[
 C_p=\{v^{b_v+1}:v\in E\},\qquad
 C_a=\{p^0\}\cup\{v^{b_v+1}:v\in E\}.
\]
Thus $r$ is oldest at both caches. The request in $a\to r$
is $\{a^0,p^1\}\cup\{v^{b_v+2}:v\in E\setminus\{r\}\}$.
All these identities are valid and protected at $r$; duplicate
ages may decrease under MIN. The frozen cache of $p$ still
permits $p\to r$, whose request is
$\{p^0\}\cup\{v^{b_v+2}:v\in E\setminus\{r\}\}$.
All requested identities are already at $r$, so no eviction
occurs there and its $a^0$ survives. The full reply restores
$a^0$ at $p$. No other cache changes.
\end{proof}

\begin{lemma}[Capacity-three empty-owner terminal gain]
\label{lem:c3ownerterminalgain}
Let $c=\ell=3$, $C_a=\varnothing$, and suppose a holder $p$
caches a nonholder and can legally select $a$. There is a
finite path increasing $N_a$, with every initial holder
present at the endpoint.
\end{lemma}
\begin{proof}
If some oldest non-$a$ entry of $p$ is a nonholder, apply
Lemma~\ref{lem:fullshuffletrianglegain}. Otherwise
$C_p=\{a^A,r^R,s^S\}$, where $r$ is a nonholder, $s$ is
a holder, and $A\ge S>R$. Execute $p\to a$, then $a\to s$,
retaining $r$ at the owner. This is legal: $s$'s reply
contains self-$a$, so at most two usable identities remain;
if a victim is needed, the other resident $p$ is eligible.
The responder now has identity set $\{a,p,r\}$, with ages
$0,e,d$ satisfying $e\le1$ and $d\le R+2$.
The frozen $p$ still has $s^{S+1},r^{R+1}$, so $p\to s$
is legal. Its request $\{p^0,r^{R+2}\}$ contains only
identities already at $s$: $s$ retains $a^0$ and its full
reply restores $a^0$ at $p$. Every original holder is now
present, while $a$ still caches the untouched nonholder $r$.
Corollary~\ref{cor:c3zeroify} gives strict holder gain.
\end{proof}

\begin{lemma}[Length-three empty-owner terminal gain at every capacity]
\label{lem:threeownerterminalgain}
Let $\ell=3$, $c\ge3$, and $C_a=\varnothing$. If a holder
$p$ can select $a$ as oldest and also caches a nonholder,
a finite legal path increases $N_a$ and preserves every
initial holder at its endpoint.
\end{lemma}
\begin{proof}
Capacity three is Lemma~\ref{lem:c3ownerterminalgain}.
Suppose $c\ge4$. Choose an oldest non-$a$ residual $s$ at
$p$, preferring a nonholder in a tie. In $p\to a$, send
$s$ and, if a second extra is required, a residual $t$.
If $s$ is a holder, choose the specified nonholder as $t$;
otherwise any second residual suffices. The empty owner
receives $p^0,s^{S+1}$ and optional $t^{T+1}$, with $S\ge T$.
Thus $a\to s$ is legal. When $s$ is a holder, require $a$
in its reply. The request installs and protects
$a,p$ and optional $t$ at $s$; MIN may lower duplicate ages.

The cache of $p$ has been frozen since its first exchange,
so $p\to s$ is still legal. Send $t$ if present, and one
additional residual $u$ if required by the fixed buffer size.
Thus the request is fresh $p$ together with optional $t,u$.
Only $u$ can be novel at $s$. Require $a,p$ in the reply.
If $s$ is full and $u$ is novel, choose the third reply entry
$x$ outside $\{a,p,t\}$ and evict it at $s$. Such an $x$
exists because $c\ge4$ and cannot equal the novel $u$.
Otherwise no eviction is needed at $s$. Its $a^0$ therefore
survives. At active $p$, self-$p$ is filtered from the reply,
leaving at most two usable identities; its preceding removals
of $a$ and $s$ leave at least two vacancies. Hence $p$ regains
$a^0$.

If $s$ was a nonholder, this is already strict gain. If $s$
was a holder, every original holder has been restored and
$t$ is still an untouched nonholder. The owner retained $t$
in $a\to s$: its reply contains self-$a$, leaving at most
two usable identities beside the two retained residents in
capacity at least four. The owner was frozen afterwards.
Corollary~\ref{cor:threegeneralzeroify} now gives strict gain
from this cached nonholder. No other initial holder is lost.
\end{proof}

\begin{lemma}[Outward holder gain at full capacity three]
\label{lem:c3outwardgain}
Let $c=\ell=3$, $C_a=\varnothing$, and suppose exactly $h$
peers hold identity $a$, all at age zero. If a holder caches
a nonholder, a finite legal path attains $N_a>h$.
\end{lemma}
\begin{proof}
Maintain a current holder $p$ caching a nonholder. Its anchor
age may become positive, but every other holder remains zero,
and the owner stays empty until a terminal gadget begins.
Choose an actual oldest partner. An owner contact terminates
by Lemma~\ref{lem:c3ownerterminalgain}.

A contact to another holder $q$ also terminates. Necessarily
$C_p=\{a^A,q^Q,r^R\}$ with nonholder $r$. The exchange
$p\to q$ refreshes $p$'s anchor, and leaves
$C_q=\{a^0,p^0,r^d\}$ with $0\le d\le R+1$.
Choose $q\to r$, an oldest choice, possibly a tie. If $q$
retains $a$, there is already a gain. If it loses $a$, the
new holder $r$ has $\{q^0,a^1,p^e\}$ with $e\le1$.
All other holders remain zero, $q$ is a nonholder, and $a$
is an oldest entry at $r$. The empty-owner terminal lemma
therefore applies to current $r$.

It remains to consider a nonholder partner $q$. If $p$ can
retain its anchor, choose that outcome and gain a holder.
Otherwise $q$ must be full and not cache $p$. Indeed, a
self-filtered reply of at most two identities permits retention
of $a$: if both are novel, sacrifice the other extra; if that
extra is received, at most one identity is novel. For an
underfull $p$ with cache $\{a,q\}$, two identities simply fit.
In a loss case the active final identity set is the old
$C_q$; duplicate numerical ages may MIN-decrease. Thus $p$
becomes a nonholder. If it had a second extra $r$, the new
current $q$ has $\{p^0,a^{A+1},r^d\}$, $d\le R+1$.
If $p$ was underfull, choose one old responder identity $D$
as survivor, giving $\{p^0,a^{A+1},D\}$. All these choices
are legal; the new current caches its nonholder predecessor,
while all other holders remain zero.

An infinite unfinished trajectory would consequently consist
only of such nonholder migrations. Its current anchor age is
exactly $A+t$ after $t$ migrations, with no duplicate refresh.
Let $M_s(t)$ be the global maximum age of identity $s$ (zero
if absent), and $K_s(t)=M_s(t)-t$. Each $K_s$ is nonincreasing.
When $s\ne a$ is selected, its active copy is removed and
not sent, and its owner cannot return self-$s$. All other
$s$ copies have unchanged or reduced ages, so
$K_s(t+1)\le K_s(t)-1$. Oldest legality requires
$K_s(t)\ge A$ at a selection. Each $s$ can therefore be
selected at most $(M_s(0)-A+1)_+$ times. Finitely many
identities cannot support infinitely many migrations.
A terminal gadget or direct gain is reached in finite time;
every prescribed finite choice has positive probability.
\end{proof}

\begin{lemma}[Outward holder gain for length three and capacity at least four]
\label{lem:threesparseoutwardgain}
Let $\ell=3$, $c\ge4$, and $C_a=\varnothing$. Suppose
exactly $h$ peers hold identity $a$, all at age zero, and
some holder caches a peer other than $a$ that lacks $a$.
A finite legal path attains $N_a>h$.
\end{lemma}
\begin{proof}
Call peers other than $a$ lacking $a$ nonholders. Maintain
a current holder $p$ with a cached nonholder target $r$.
Let $Z$ be the initial holders other than the initial current.
Until termination, only the current initiates; all members of
$Z$ remain passive zero holders, and the holder set is exactly
$Z\cup\{p\}$. A migration changes $p$ and makes its previous
initiator the new designated target. Define the fixed set
$Y=V\setminus(Z\cup\{a\})$ and
\[
 F(x)=\sum_v|\operatorname{ids}(C_v)\cap Y|,\qquad
 \Phi=(F,c-|C_p|).
\]
Thus $Y$ consists of the current and the nonholders. The
lexicographic potential takes finitely many values. All counts
below are identity counts, unaffected by numerical MIN merges.
An owner contact terminates by
Lemma~\ref{lem:threeownerterminalgain}.

Consider a zero-holder partner $q\in Z$. Write $C,D$ for the
pre-exchange current and responder identity sets, and $d=|C|$.
Always require $a$ in the reply, preserving the passive zero
and refreshing the current anchor.
If $q$ knows both $p,r$, send $a,r$ and return $a,p,r$.
All received identities are duplicates or self-filtered, so
$F$ does not decrease and the current cache shrinks by one.
If $q$ knows $p$ but not $r$, send $a,r$ and return $a,p$ plus
a residual $D_0$ if required. A full responder can choose
$D_0\notin C$ as its victim: its $c-2$ residents outside
$\{a,p\}$ overlap $C$ in at most $d-3\le c-3$ identities.
Any counted victim is copied to a new active location, while
the counted target $r$ gains a new passive copy. At most one
active novel identity arrives, so no active victim is needed.
An underfull responder needs no victim. In either case $F$
strictly increases.

If $q$ does not know $p$, send $a,B$. For a full current,
choose $B\ne r$, possible since $c\ge4$, leaving $r$ unsent.
For an underfull current any second extra is allowed: after
partner deletion two novel reply identities fit, so $r$ still
cannot be evicted. If $B\in D$, \emph{require $B$ in the reply}
beside $a$. Only fresh $p$ is novel passively. If a victim is
needed, choose the third reply identity outside $C$; it exists
because $|D\setminus C|\ge |D|-d+1\ge1$ for a full responder.
That victim is copied to a new active location. Returned $B$
is protected actively, and at most one other novel identity
arrives. Fresh $p\in Y$ is an additional new passive entry,
so $F$ strictly increases. Returning $B$ is essential here.

If $B\notin D$, both $p,B$ are novel passively. With $e=|D|$,
\[
 |D\setminus C|\ge e-d+2\ge e-c+2,
\]
since $q,B\in C$ are absent from $D$. Choose all
$(e+2-c)_+$ needed passive victims among these absent-current
residents and include them with $a$ in the reply; at most two
are needed. Every counted passive victim is copied into a new
active location. The only possible active extra victim is $B$;
if lost, it has just been copied into a new passive location.
Fresh $p$ is an additional counted entry. Thus $F$ strictly
increases, also for shorter underfull replies. In all these
zero-holder cases the target remains cached and $\Phi$
strictly increases.

Now let $q$ be a nonholder. Include $a$ in the request. If
the current is underfull, it can retain $a$: at most three
novel replies and at least two vacancies require at most one
victim, which can be the other sent extra $B$. If $B$ is
returned it is a duplicate, so no victim is needed. A current
of size two needs no extra $B$, since its sole retained anchor
plus three replies fit in $c\ge4$. This gives direct gain.
For a full current, if $q$ knows $p$, return $p$, leaving at
most two usable identities and the same retention choice.
If instead $q$ knows a current residual $x\notin\{a,q\}$,
return $x$ and send a different residual $B\ne x$ with $a$;
there are at least two residual choices. If $B$ is also
returned there is no overflow; otherwise evict $B$, not $a$.
This again gives gain.

Consequently an unfinished migration need only occur from a
full current to a responder lacking $p$ and disjoint from $C$.
A reply shorter than three permits retention of $a$, so assume
three novel reply identities. Active victims are exactly $a,B$.
The responder receives three novel protected identities $p,a,B$;
choose every necessary passive victim among its three replied
residents. Each counted passive victim is copied to a new active
location, and lost counted $B$ is newly copied passively. Fresh
$p\in Y$ compensates the removed selected $q\in Y$, giving
$F'\ge F$. The responder becomes current, with fresh predecessor
$p^0$ as designated nonholder target. Its anchor age is exactly
the old current age plus one. All of $Z$ stays zero. The old
current was full, so $c-|C_p|=0$; the new vacancy count is
nonnegative, and $\Phi$ cannot decrease even if $F'=F$.

An infinite unfinished trajectory would eventually have constant
$\Phi$, and hence no zero-holder contacts. Its tail consists
only of migrations. Starting that tail with anchor age $A$,
the current anchor has age $A+t$ after $t$ events. For every
$s\ne a$, put $K_s(t)=M_s(t)-t$, where $M_s$ is its global
maximum stored age, zero if absent. Each $K_s$ is nonincreasing;
when $s$ is selected, its active copy is removed and not sent,
its owner cannot return self-$s$, and all other copies have
unchanged or reduced ages. Hence $K_s$ drops by at least one.
Oldest legality requires $K_s(t)\ge A$ at a selection, allowing
at most $(M_s(0)-A+1)_+$ selections of each identity. This
contradicts an infinite tail. Thus a finite legal path reaches
direct gain or the terminal gadget, with positive probability.
\end{proof}

\begin{theorem}[Uniform sampling at full capacity three]
\label{thm:c3generaluniform}
For every $n\ge4$, the adopted kernel with $c=\ell=3$
converges from every weakly connected initial state with
arbitrary finite nonnegative integer ages to a unique
stationary law $\pi$. There are finite $T(n)$ and
$\gamma(n)>0$ such that
\[
 \inf_x\widehat P_3^{T(n)}(x,\{H(V,3)\})\ge\gamma(n),\qquad
 \sup_x\|\widehat P_3^t(x,\cdot)-\pi\|_{\TV}
 \le(1-\gamma(n))^{\lfloor t/T(n)\rfloor}.
\]
Every observer and request size $1\le b\le3$ has positive
stationary availability and a uniform limiting ordered sample
conditional on availability. No polynomial or practical bound
on these constants is asserted. This theorem does not cover
sparse $c\ge4$ at shuffle length three.
\end{theorem}
\begin{proof}
Fix a closed communicating class and identity $a$. By
Corollary~\ref{cor:c3zeroify}, if $k=\max Z_a$ in this class,
then $N_a\le k$ throughout it, and all holders at a maximizer
are zero. Also $k\ge1$: if no peer caches $a$, weak
connectivity forces the owner nonempty and an owner activation
creates a zero. Write $I$ for a maximizing zero-holder set.

There is no directed cache path from $a$ to a nonholder.
Otherwise take a simple such path with internal vertices in
$I$. Apply Lemma~\ref{lem:c3learnedcallback} with $H$ equal
to $I$ minus its remaining path vertices, retaining the first
path head. Before the first boundary contact those path rows
are untouched. A nonholder contact contradicts $N_a\le k$;
otherwise the contacted path vertex returns its next head,
which is received and protected at the owner. Repeat with
the shorter suffix. This eventually gives the contradiction.
No prescribed-target service across another nonholder is used.

Let $R$ be the set directed-reachable from $a$, including $a$.
Then $R\setminus\{a\}\subseteq I$ and its cache identities
stay in $R$. Freeze its complement. Inspect an actual oldest
owner head $Q$, and reserve $a^0$ while normalizing $Q$'s
residual weak component in $R\setminus\{a\}$ with root $Q$.
The effective capacity and length are both $\min(2,m-1)$
for component size $m$; a singleton needs no exchanges.
Lemma~\ref{lem:anchorreduction} and the lower-length theorems
apply, omitting empty residual activations. Now $C_Q=\{a^0\}$.
The frozen owner's still-legal contact to $Q$ removes its head
without replacement; the responder's fresh anchor and at most
two extras fit. Recompute components and repeat. This produces
an empty-owner maximizer without presupposing outward closure.

At every empty-owner maximizer, Lemma~\ref{lem:c3outwardgain}
forbids an edge from a holder to a nonholder. Among these
states choose one minimizing the number $d$ of residual weak
components on $I$. If $J=V\setminus(I\cup\{a\})$ is nonempty,
weak connectivity supplies an edge from $J$ to $I$: the owner
is empty, no nonholder caches $a$, and there is no outward
holder edge. Freeze $I$ and $a$, and apply
Theorem~\ref{thm:threefirstcontact} until an outside initiator
$P$ is about to contact $Q\in I$. No outside peer has yet
acquired $a$. Pause the outside and normalize $Q$'s residual
component with root $Q$. The owner stays empty, $d$ is unchanged,
and $P$'s pending oldest choice is unchanged.

Now $C_Q=\{a^0\}$. If $P$ has at most two entries, its request
fits beside that anchor and its reply creates $a^0$ at $P$,
a contradiction. Otherwise $C_P=\{Q,R,S\}$, and the endpoint
identity sets are
\[
 C_P'=\{a^0,R^u,S^v\},\qquad C_Q'=\{P^0,R^u,S^v\}.
\]
The owner's cache is still empty and the new zero set is
$I'=(I\setminus\{Q\})\cup\{P\}$. Both $R,S$ must lie in
$I\setminus\{Q\}$, or $P$ has a forbidden outward edge at
this new empty-owner maximizer. If $Q$'s normalized component
was nontrivial, some remaining holder still caches its root
$Q$, another forbidden outward edge. Thus it was a singleton.
Removing that component and attaching $P$ to $R,S$ in the
remaining components strictly decreases $d$, a contradiction.
Therefore $I=V\setminus\{a\}$.

Every closed class reaches a universal zero anchor, hence
any prescribed labelled $H(V,3)$ by
Proposition~\ref{prop:threezeroanchor}. For completeness,
closed-class reachability is not an additional assumption.
Theorem~\ref{thm:thesisgeneralreset} supplies a finite reset
family $F$, horizon $B$, and probability $\delta>0$.
In the finite reachability graph on $F$, every vertex reaches
a terminal strongly connected component. Every state reachable
from such a component has a finite path back to $F$ by reset;
terminality puts that return in the same component. Those
reachable states thus form a closed communicating class.
Consequently every member of $F$ has a finite path to the
same prescribed atom. Choose one path per member and take
their maximum length $L$. Pad each path to length $L$ at
the atom's empty-root self-loop, and let $\varepsilon>0$
be the minimum probability of these padded paths.
Then $\widehat P_3^{B+L}(x,\{H(V,3)\})\ge\delta\varepsilon$
for every $x$. Lemma~\ref{lem:atom} proves the stated uniform
convergence. Relabelling gives positive stationary mass to
an atom where any specified observer has three entries;
Theorem~\ref{thm:stationary} gives conditional uniformity.
\end{proof}

\begin{theorem}[Uniform sampling for every three-entry shuffle]
\label{thm:threegeneraluniform}
For every $n\ge4$ and $3\le c\le n-1$, the adopted kernel
with fixed shuffle length $\ell=3$ converges from every weakly
connected initial state with arbitrary finite nonnegative integer
ages to a unique stationary probability law $\pi$. There are
finite $T(n,c)$ and $\gamma(n,c)>0$ such that
\[
 \inf_x\widehat P_3^{T(n,c)}(x,\{H(V,c)\})\ge\gamma(n,c),\qquad
 \sup_x\|\widehat P_3^t(x,\cdot)-\pi\|_{\TV}
 \le(1-\gamma(n,c))^{\lfloor t/T(n,c)\rfloor}.
\]
For every observer and $1\le b\le c$, stationary availability
is positive and the limiting ordered sample conditional on
availability is uniform over distinct nonself labels. No
polynomial or practical mixing bound is asserted, and shuffle
lengths at least four are not covered by this theorem.
\end{theorem}
\begin{proof}
Capacity three is Theorem~\ref{thm:c3generaluniform}; assume
$c\ge4$. In a closed communicating class fix identity $a$
and maximize $Z_a=k$. Corollary~\ref{cor:threegeneralzeroify}
gives $N_a\le k$ throughout the class, and all holders at a
maximizer form a zero set $I$. Also $k\ge1$, since in the
absence of any copy weak connectivity forces a nonempty owner
cache, whose activation creates a zero.

No nonholder is directed-reachable from the owner. Otherwise
take a simple path with internal vertices in $I$, retain its
first head, and apply Lemma~\ref{lem:threegeneralcallback}
with $H$ equal to $I$ minus the remaining path vertices.
Before its terminal contact all those path rows are untouched.
A nonholder contact gives forbidden holder gain. At a path
vertex require the reply to contain $a$ and the next path head,
adding an arbitrary third entry if needed. These two entries
fit the fixed three-entry reply; the callback permits this
choice at its final boundary event. The next head is received
and protected at the owner. Repetition along strictly shorter
suffixes reaches a nonholder, a contradiction. This explicitly
uses a chosen partial reply, not an assumed full-cache reply.

Let $R$ be the owner's directed reachable set, including $a$.
Then $R\setminus\{a\}\subseteq I$ and every row in $R$
contains only identities from $R$. Freeze the complement.
For an actual oldest owner head $Q$, reserve $a^0$ and
normalize its residual weak component with root $Q$. On a
component of size $m$, the effective capacity and length are
$\min(c-1,m-1)$ and $\min(2,m-1)$. The lower-length theorems
and Lemma~\ref{lem:anchorreduction} apply, with singleton
components requiring no action and empty residual no-ops
omitted. The owner was frozen, so its contact to the prepared
$C_Q=\{a^0\}$ is still legal. The sole self-$a$ reply removes
one owner head without replacement. At the responder fresh
$a$ and at most two extras fit, preserving its zero. Recompute
residual components and repeat until the owner is empty.
Every original zero is preserved. Thus an empty-owner maximizer
exists without assuming global outward closure.

At any such maximizer, Lemma~\ref{lem:threesparseoutwardgain}
rules out a holder-to-nonholder edge. If
$J=V\setminus(I\cup\{a\})$ were nonempty, weak connectivity
would therefore supply an edge from $J$ to $I$: the owner is
empty and no nonholder caches $a$. Freeze $I$ and $a$, and use
Theorem~\ref{thm:threefirstcontact} to pause before the first
outside-to-$I$ exchange $P\to Q$. No outside peer can have
acquired $a$ before that crossing. Freeze the outside and
normalize $Q$'s residual component with root $Q$, using the
same effective lower-length parameters. This leaves the owner
empty and $P$'s pending oldest choice unchanged.

Now $C_Q=\{a^0\}$. Resume $P\to Q$. Its sole reply supplies
$a^0$ to $P$, with space provided by mandatory partner removal.
The request contains at most three identities, none equal to
$a$, and fits beside $Q$'s old anchor in capacity $c\ge4$.
Thus $Q$ keeps its zero and $P$ becomes an additional zero
holder, contradicting maximality. Every closed class hence
contains a universal zero anchor and reaches any prescribed
labelled $H(V,c)$ by Proposition~\ref{prop:threezeroanchor}.

To make the stochastic conclusion explicit, use the finite
reset family $F$ and uniform reset horizon $B$ and probability
$\delta$ from Theorem~\ref{thm:thesisgeneralreset}. In the
finite directed reachability graph on $F$, each vertex reaches
a terminal strongly connected component. Any state reachable
from such a component returns to some member of $F$ along a
finite positive path by reset; terminality puts that member
in the same component. All states reachable from the component
therefore communicate with it and form a closed class. Thus
every reset state has a finite path to the prescribed atom.
Choose one per reset state, pad to a common finite length $L$
at the atom's empty-root self-loop, and let $\varepsilon>0$
be the minimum probability of the padded paths. Then
\[
 \widehat P_3^{B+L}(x,\{H(V,c)\})\ge\delta\varepsilon
 \quad\hbox{for every }x.
\]
Lemma~\ref{lem:atom} gives the stated uniform geometric bound.
Relabelling gives positive stationary mass to an atom where
any specified observer has a full cache, and
Theorem~\ref{thm:stationary} gives conditional uniformity for
all $b\le c$. The finite age-dependent accessibility paths
have been made uniform only through the separate reset; no
efficient bound on their lengths or probabilities is claimed.
\end{proof}

\section{Structural induction and the independent full-capacity theorem}
Write $C_L$ for the uniform canonical common-atom conclusion at length
$L$, $F_L$ for the holder-only uniform first-contact conclusion across
a frozen cut, and $U_L$ for finite-path normalization of a universal
zero anchor. These are proof-interface names, not extra assumptions.
The following proposition obtains $F_L,U_L$ from the earlier-length
inputs. The full-capacity theorem then closes its own double induction
before it is used in the arbitrary-capacity source argument.


\begin{proposition}[Structural induction from lower shuffle lengths]
\label{prop:shuffleinductionstructure}
Fix $L\ge3$. Assume uniform canonical common-atom minorization
at every smaller length $1\le j<L$, for every allowed population
and capacity, and assume the uniform holder-only frozen-cut
first-contact conclusion at length $L-1$. Then at length $L$,
for every $n\ge L+1$ and $L\le c\le n-1$:
\begin{enumerate}
\item the holder-only first-contact conclusion of
Theorem~\ref{thm:threefirstcontact} holds with a finite horizon
and positive probability bound independent of initial ages;
\item a universal zero anchor has a finite legal path to any
prescribed labelled canonical atom $H(V,c)$.
\end{enumerate}
These are structural conclusions, not convergence from arbitrary
states at length $L$.
\end{proposition}
\begin{proof}
Repeat the minimum-age-holder construction of
Theorem~\ref{thm:threefirstcontact}. Reserving target $r$ at
its minimum-age holders gives the exact length-$L-1$,
capacity-$(c-1)$ residual kernel. A residual oldest choice
is legal unless its maximum age is below the reserved anchor;
in that case contact with owner $r$ already succeeds.
Open residual cuts use the assumed lower first-contact result
on $n-1$ identities. Closed components of size $m\ge2$ have
\[
 d_m=\min(c-1,m-1),\qquad e_m=\min(L-1,m-1)<L.
\]
The assumed common atom at length $e_m$ applies, including
clipped singleton and two-entry cases; a one-vertex component
contacts the anchor owner immediately. Thus a finite maximum
$E$ of the lower first-contact and component-atom horizons,
plus one event, removes a minimum copy or contacts the cut.
At most $s$ phases raise the outside minimum by one.

Put $R_s=(sc+1)c$ and raise that minimum $R_s+1$ times.
The final old-anchor/young-count routing proof transfers
without a new induction assumption: Lemma~\ref{lem:generalcolor}
is valid at every length, and reserving the old anchor plus
a required young difference needs only two extra slots,
available for $L\ge3$. Its reply bound is
$|D_q|+|Q\cap U|\le\min(L,|Q|)$. The same bounded pair
$(sc-F,c-1-M)$ decreases at each nonterminal step, so this
last stage takes at most $R_s$ events. In particular one may
use
\[
 T_L=sE(R_s+1)+R_s,\qquad
 \eta_L=(nc\,16^c)^{-T_L}.
\]
Only target holders initiate along the controlled path;
existing responder copies are returned and protected. Before
any cut crossing, arbitrary internal events preserve some
outside target copy, so the uniform block bound also gives
the same geometric first-crossing tail as before. These
constants are finite, not claimed polynomial or practical.

For a universal zero anchor $a$, reserve it and normalize
residual components with the displayed effective parameters.
Freeze the owner while normalizing its inspected oldest head
$Q$ to a root with $C_Q=\{a^0\}$. The owner then deletes that
head without replacement. Fresh $a$ and at most $L-1$ extras
fit at $Q$ because $c\ge L$. Repetition empties the owner.
The four-exchange component merges and root changes of
Proposition~\ref{prop:threezeroanchor} use caches of size at
most two, so are unchanged. A final residual normalization
produces the prescribed labelled atom. Shorter kernels were
used only through exact reserved-anchor simulations, not
through shorter original shuffles.
\end{proof}

\begin{theorem}[Every shuffle length at full cache capacity]
\label{thm:fullcapacityalllengths}
For every $n\ge2$, let $c=n-1$ and $1\le\ell\le n-1$.
The adopted mandatory-removal, protected-MIN protocol with iid
uniform global initiators has a uniform finite-horizon common atom
$H(V,n-1)$ over all weakly connected finite-integer-age starts.
It has a unique stationary full-state law and converges to it.
Every allowed request size has positive stationary availability,
and its deterministic-time conditional local sample converges to
uniform ordered sampling without replacement.
The finite minorization constants are implicit; no polynomial
mixing bound is asserted. This theorem does not cover $c<n-1$.
\end{theorem}

\begin{proof}
Induct first on $L=\ell$, starting from the established lengths
one through three, and within a new length $L\ge4$ on
$n=L+1,L+2,\ldots$. The outer hypothesis is only the full-capacity common-atom
theorem at every smaller length and every allowed population.
The inner hypothesis is the full-capacity common-atom theorem at length $L$
for smaller populations. No sparse-capacity theorem at an
unproved length is assumed.

The first-contact and universal-zero-anchor constructions in
Proposition~\ref{prop:shuffleinductionstructure} remain inside
this family. Reserving an anchor leaves capacity $n-2$ on
$n-1$ identities; a residual component of size $s$ has effective
capacity $s-1$ and length $\min(L-1,s-1)<L$.
Thus all component atoms use only the outer hypothesis.
The open-cut part uses the previously established full-capacity
$F_{L-1}$. Starting with the length-three contact theorem, construct
$F_L$ and $U_L$ before the population induction for length $L$.
If needed, a connected auxiliary cut completion is obtained by
filling one frozen outside row with all other reduced identities;
no preboundary exchange uses that row. There is no $F_L/C_L$
circularity.

Fix a closed communicating class $\mathcal C$.
At capacity $n-1$ no merge ever evicts: there are only $n-1$
possible nonself identities. The underlying undirected graph
therefore only gains edges. Mandatory partner removal preserves
its edge through the fresh reverse descriptor at the responder.
Communication within $\mathcal C$ makes this one fixed connected
graph $G$ throughout the class.

\emph{A finite frozen-source erasure.}
Choose a peer $q$ and a state attaining its maximum row size
$m_q$ over $\mathcal C$, and freeze $q$. If $m_q<n-1$, a novel
passive insertion would exceed that maximum, so its identities
cannot change. If $m_q=n-1$, its row already contains all possible
identities. In either case its numerical ages can only decrease
by MIN. Let $M$ bound its ages in this initial state.

Let $R$ be the old-lineage elimination horizon from the general
colored-reset proof, for example
\[
 R=(nc+1)(nc^2+1)(ncL^c+1),\qquad B=M+(nc+1)R.
\]
At any later state, one of the $nc+1$ intervals
$(M+jR,M+(j+1)R]$, $0\le j\le nc$, is empty of stored ages.
Color ages at most its lower endpoint young and those above
its upper endpoint old. The unpadded reset path initiates only
old-holding peers and eliminates all old descriptors within $R$
events. Peer $q$ stays young, since its row remains bounded by
$M$, and never initiates. This gives a $q$-frozen reset to ages
at most the same fixed $B$.

While $q$ is frozen, the global minimum stored $q$ age cannot
decrease: no fresh $q$ is generated. If copies remain, let their
minimum be $k$, and let $S$ be its holders. The minimum-holder
phase in the structural first-contact argument applies with
$q$ reserved. Internal $S$ contacts preserve both $k$ copies by
$\min(k+1,k)=k$. An open residual cut is reached using $F_{L-1}$;
a closed residual component is normalized using the outer
common-atom hypothesis. If a proposed residual oldest is younger
than $k$, contact $q$ directly instead; ties permit the residual
step. A singleton or normalized empty residual row also contacts
$q$. All initiators belong to $S$, hence differ from $q$.
At a boundary the outgoing $q$ age is $k+1$, and the outside
responder has no $q$ copy or an older one. At a direct owner
contact mandatory deletion removes the active $q$ copy.
Either terminal loses a $k$ copy and creates none. Repeating
at most $n-1$ such finite phases raises the minimum or erases
$q$ entirely.

Repeat until $q$ is absent or its minimum exceeds $B$, then
apply the $q$-frozen reset. A surviving $q$ copy could not have
age at most $B$, so all incoming $q$ references are erased.
This is a finite argument with $M,B$ fixed throughout, not a
claim about an infinite prescribed schedule or a contact-rate
extinction estimate. Every history used stays in $\mathcal C$.
The frozen row of $q$ now represents all its incident edges;
its identity set is exactly $N_G(q)$ and is nonempty.

\emph{Connected complement.}
Put $W=V\setminus\{q\}$. If $G[W]$ had multiple components,
let $u$ be an oldest head of $q$ and choose a cached neighbor
$v$ in another component. Such neighbors exist in every
component by connectivity of $G$. Since $L\ge2$, a legal
$q\to u$ request can include $v$ as an extra. Protection then
creates edge $uv$, contradicting the fixed graph $G$.
Thus $G[W]$ is connected. This contradiction need not be executed;
retain the frozen-source state just constructed.

\emph{One source contact produces a universal zero.}
No row in $W$ contains $q$. Its $W$-only transition rule is exactly
that of capacity $n-2$ and length $\min(L,n-2)$ on $n-1$ peers.
This is a conditional-on-initiator or finite controlled-path
simulation: actual prescribed initiators still have probability
$1/n$. At $n=L+1$ the reduced length is $L-1$, covered by the
outer induction; otherwise the same length on a smaller population
is covered by the inner induction.

Inspect an oldest head $u$ of frozen $q$. Normalize $W$ to its
canonical atom with root $r\ne u$, possible since $|W|\ge2$.
All peers of $W\setminus\{r\}$ then cache $r^0$.
The still-legal contact $q\to u$ may require $r^0$ in the reply.
Peer $q$ receives it, and $u$ retains its old $r^0$ because there
are no evictions and MIN cannot increase it. All other copies
are unchanged. Thus $r^0$ is universal. The already established
$U_L$ construction reaches the prescribed labelled $H(V,n-1)$.
Every closed class reaches this same atom.

Finally apply the independent uniform finite-age reset.
Its finite reset-family reachability graph has terminal strongly
connected components generating genuine closed communicating
classes: another reset from any reachable state returns to that
finite family, and terminality forces return into the same
component. Each reset state therefore has a finite path to the
common atom. Choose these finitely many paths, pad at the atom's
empty root to a common length, and only then take their minimum
positive probability. Composing with the initial reset gives
uniform minorization, without presuming irreducibility or
aperiodicity. Lemma~\ref{lem:atom}, label symmetry and relabelled
full-cache atoms give convergence, positive availability and
uniform conditional samples. This closes both inductions.
\end{proof}

\section{The general arbitrary-capacity proof}

\begin{lemma}[Partial colored reset with one excluded initiator]
\label{lem:partialcolorsource}
Fix $2\le L\le c\le n-1$ and a peer $q$. Color descriptors old or
young as in Lemma~\ref{lem:generalcolor}, assume every descriptor of
identity $q$ is old, and never initiate $q$. Suppose the old-oldest and
mixed-merge color assumptions of that lemma remain valid throughout the
following horizon. There is a controlled path erasing every $q$
descriptor within $H$ events, where one conservative choice is
\[
 K=n(c+1)^c,\quad D=ncL^c,\quad
 R=(nc+1)(nc^2+1),\quad
 H=R\bigl((K+1)(D+1)+1\bigr).
\]
Other old descriptors need not be erased. In particular the frozen
row $C_q$ is allowed to contain old descriptors throughout.
\end{lemma}
\begin{proof}
Stop when $q$ has no copies. Otherwise choose $p\ne q$ with an old
descriptor maximizing $f_p$ among such peers, and choose an actual
old oldest responder $r$. Use exactly the extras, joint reply/victim
choice, and active old-first victims of Lemma~\ref{lem:generalcolor}.
Write $F=\sum_i f_i$ and $V=\sum_i f_i^2$.

The proof that $F$ cannot decrease does not use maximality. If
$F'=F$, that proof gives $p\in Y_r$ and $Y_p\subseteq Y_r$, hence
$f_r\ge f_p+1$. There is no young loss at $p$, and the only young
changes are $t$ responder victims transferred to new young locations
at $p$. The responder victim exclusions protect $Y_p\cup\{p\}$,
so $0\le t\le f_r-f_p-1$. Therefore
\[
 V'-V=2t(t-(f_r-f_p))\le0,
\]
with equality only for $t=0$. This calculation remains valid when
$r=q$ has old entries: old arrivals do not change the young counts.
On a consecutive plateau of both $F$ and $V$, every young identity
set, and hence every $f_i$, is fixed.

On such a plateau, any responder $r\ne q$ is entirely young before
the event. Otherwise $f_r>f_p$ contradicts the maximal choice of $p$.
More strongly, no outside old-holder can be a responder at all.
Only a contact with $q$ can return old descriptors, and that reply
cannot contain identity $q$.

Call a tenure the persistence of one $q$ copy at an outside row,
starting at the beginning of the plateau or at its insertion, and
ending when it is lost or the plateau ends. Its row is never passively
updated during the tenure, since it is an outside old-holder. Every
non-$q$ contact initiated there removes an old head and receives only
young descriptors. A contact with $q$ removes the $q$ copy itself.
Consequently a tenure has at most $c$ initiations. A new $q$ copy
can arise only as a request extra from an existing tenure, and its
responder has strictly larger fixed $f$. Each initiation creates at
most one such child tenure. Since an old-holder has $f\le c-1$,
these genealogies have depth at most $c-1$, branching at most $c$,
and at most $n-1$ roots. Their total number, and thus the number of
contacts with $q$ on the plateau, is at most
\[
 (n-1)\sum_{j=0}^{c-1}c^j\le K.
\]
Fresh $q$ cannot appear, because $q$ never initiates; thus all its
copies remain old and this accounts for every insertion.

At every plateau event not contacting $q$, the responder is old-free,
so the original strict descent of
$W=\sum_i o_iL^{c-f_i}$ applies. Its range is $0\le W\le D$.
At most $K$ contacts with $q$ interrupt this descent. Hence a plateau
with surviving $q$ copies has at most $(K+1)(D+1)$ internal events.
There are at most $R$ successive $(F,V)$ plateaus, because $F$ can
only increase, and $V$ decreases while $F$ stays fixed. The stated
$H$ also accounts for events moving between plateaus. Thus the
controller must erase $q$ before exhausting the horizon.
\end{proof}

\begin{lemma}[Source creation from lower-length inputs]
\label{lem:generalsourcecreation}
Fix $L\ge4$ and suppose the smaller-length common-atom and
length-$(L-1)$ first-contact hypotheses of
Proposition~\ref{prop:shuffleinductionstructure} hold. From every
weakly connected state with finite nonnegative integer ages, and
for any specified peer $q$, a finite legal path with no initiation
of $q$ reaches a state where $q$ has no incoming references and
$C_q$ is nonempty.
\end{lemma}
\begin{proof}
If $q$ is already absent, weak connectivity implies $C_q\ne\varnothing$.
Otherwise its minimum stored age cannot decrease while $q$ does
not initiate. Use the minimum-holder phase from the structural
first-contact proof. At current minimum $k$, reserve $q$ in the
minimum holders; their internal residual kernel has length $L-1$
and capacity $c-1$. If a proposed residual oldest has age below
$k$, contact $q$ instead; otherwise the residual step is an actual
oldest step (ties are legal). A closed residual component is
normalized using the smaller-length common atom. For an open
component, apply the lower-length first-contact theorem on the
ambient $V\setminus\{q\}$, with capacity $c-1$, length $L-1$,
and the minimum-holder set as the active set. Assign auxiliary empty
rows outside that set; the first-contact conclusion needs a cached
cut edge, not global connectivity of the auxiliary state. Outside
rows are not read before the first boundary event. At that event
use the actual legal reply, whose $q$ copy, if any, is older than
$k$, because the responder is outside the minimum-holder set.

An internal contact preserves both minimum copies through
$\min(k+1,k)=k$. A boundary contact sends $q$ at age $k+1$;
a direct owner contact deletes it by mandatory removal. Either
terminal removes a minimum copy and creates none. In at most
$n-1$ such finite phases the minimum rises or $q$ disappears.
Closed singleton residual components cause a direct owner contact. These
steps never require a bound on the passive row of $q$.

Let $H$ be the constant of Lemma~\ref{lem:partialcolorsource} and
$B=(nc+1)H$. Repeat the preceding finite phases until $q$ disappears
or its minimum exceeds $B$. Integer ages ensure that finitely many
raises suffice. If copies remain, among the $nc+1$ intervals
$(jH,(j+1)H]$, $0\le j\le nc$, one contains no current descriptor
age, since at most $nc$ descriptors exist. Put $a=jH$ for this
interval and color ages at most $a$ young, and those above $a+H$
old. All $q$ copies are old; the row of $q$ may have either color.
During the next $H$ events young ages are at most $a+t\le a+H$,
and old ages remain greater than $a+H$. Thus mixed MIN merges
keep young and every old-holder has an old actual oldest. The
partial colored reset applies and erases $q$ without initiating it.

If a last $q$ copy is lost, its final disappearance need not be a
contact with $q$: an eligible eviction can also remove it. However
$q$ has never initiated and its row cannot become empty passively:
a passive exchange receives the protected fresh initiator. If it
was initially empty and never contacted, weak connectivity would
preclude absence of $q$ at the endpoint. Hence the resulting source
row is nonempty in every case. All histories preserve weak connectivity.
\end{proof}

\begin{lemma}[Source complement joining and extension]
\label{lem:generalsourceextension}
Fix $L\ge4$. Assume the lower-length structural hypotheses above,
and common atoms at length $L$ for smaller populations. A weakly
connected state with an incoming-free peer $q$ and nonempty row
has a finite legal path to a universal zero anchor.
\end{lemma}
\begin{proof}
Set $S=C_q$ as an identity set and $W=V\setminus\{q\}$.
Every weak component $A$ of $W$ meets $S$. Normalize each component
while $q$ is frozen, ordering all vertices of $S\cap A$ above the
other vertices. A singleton is already normalized. The effective
component parameters are $\min(c,|A|-1)$ and $\min(L,|A|-1)$:
either smaller length or smaller population supplies the theorem.
The vertices of $S\cap A$ then reference only higher vertices of
that set, while every vertex of $A\setminus S$ retains its
component-maximum root in $S$ at age zero.

The set $K_0=\{q\}\cup S$ is now a closed subsystem, weakly
connected through $q$. Its effective capacity is $s=|S|$ and its
length is $\min(L,s)$, so it has full capacity. Apply the already
independent Theorem~\ref{thm:fullcapacityalllengths} to normalize
$K_0$ with $q$ least in the order. Outside rows stay frozen, even
if they have incoming references to $K_0$. At the endpoint $q$
again has no incoming references and holds precisely $S$ at zero.
The induced graph on $S$ is connected; frozen outside rows still
point to their component roots in $S$. Thus $W$ is connected.
If $K_0=V$, then $s=c=n-1$, so the full-capacity theorem remains
an independent input, not a circular sparse induction.

Normalize the connected $W$ with effective capacity $\min(c,n-2)$
and length $\min(L,n-2)$, keeping $q$ frozen. Before doing so choose
an oldest $u$ of $q$ and, if $s\ge2$, an oldest $v$ among the
remaining entries, and put $u$ greatest and $v$ second greatest
in the canonical order on $W$. For $s\ge2$, we have
$C_u=\varnothing$ and $C_v=\{u^0\}$. The exchanges $q\to u$, $q\to v$ create a
universal $u^0$: the first removes $u$ and receives no reply; the
second is still oldest, receives $u^0$, and sends at most
$\min(L,s-1)\le c-1$ identities, none equal to $u$. They fit beside
$v$'s $u^0$, so that zero is not evicted. Every other nonowner in
$W$ retains its canonical $u^0$. If $s=1$, the two exchanges
$q\to u$, $u\to q$ instead leave $q=\{u^0\}$ and $u$ empty,
with all other canonical zeros unchanged. This includes the
singleton-complement case without a zero-capacity invocation.
\end{proof}


\begin{proof}[Proof of Theorem~\ref{thm:allparametersuniform}]
Use outer induction on length, with the established lengths one,
two and three as bases, carrying the first-contact conclusion along
with the common-atom theorem. At length $L\ge4$, lower-length
common atoms and $F_{L-1}$ supply $F_L$ and the universal-zero
normalization $U_L$ by
Proposition~\ref{prop:shuffleinductionstructure}. Induct internally
on population $n\ge L+1$, for every capacity at that population.
The smallest population is full capacity and is already covered by
Theorem~\ref{thm:fullcapacityalllengths}.

Lemma~\ref{lem:generalsourcecreation} gives a source from any state.
Lemma~\ref{lem:generalsourceextension} uses only smaller-population
or smaller-length component theorems and the independent
full-capacity theorem to reach a universal zero. Then $U_L$
reaches the same prescribed labelled canonical atom. In particular
every closed communicating class reaches that atom. This reasoning
uses finite legal paths; it does not infer a uniform duration from
age-dependent minimum-holder phases.

For uniformity, apply the independent finite-age reset of
Theorem~\ref{thm:thesisgeneralreset} to a fixed finite set $\mathcal F$.
The directed reachability graph on $\mathcal F$ has terminal strongly
connected components. Each generates a closed communicating class:
any reachable state admits another reset into $\mathcal F$, and
terminality forces its return into that same component. Thus every
reset state has a finite path to a closed class and then to the
canonical atom. Choose one such path per reset state, pad them to
a common length using the atom's empty-root no-op, and take the
minimum of their finitely many positive probabilities. Composing
with the reset gives uniform finite-horizon minorization.
Lemma~\ref{lem:atom} gives full-state uniqueness and convergence.
Relabeling invariance gives conditional uniform ordered samples;
relabelled canonical atoms with the observer's row of size $c$
give positive availability. This closes the population induction,
and $F_L$ supplies the first-contact input for the next length.
\end{proof}

\end{document}
